\chapter{Flattening Stratification of a Coherent Sheaf}
\label{chap:Picard_FlatteningStratification}

\section{Setup and motivation}
\label{sec:flatstrat_setup}

This chapter establishes the existence of the \emph{flattening stratification} of
a coherent sheaf \(\F\) on a projective \(S\)-scheme: a finite collection of
locally-closed subschemes of \(S\) on each of which the restriction of \(\F\)
becomes flat, together with the universal property that a base change makes
\(\F\) flat exactly when it factors through the stratification.

The statement is a theorem of commutative algebra and scheme theory
independent of any moduli construction.  It is the technical input by which
Nitsure's exposition of Grothendieck's existence theorem
(\cite[\S5]{nitsure-hilbert-quot}) cuts the locus of \emph{quotients with flat
fibres of the prescribed Hilbert polynomial} out of a Grassmannian, once the
functor of points of \(\Quot_{\F/X/S}^{\Phi,L}\) has been embedded there via
Castelnuovo--Mumford regularity.

We work throughout with a noetherian base scheme \(S\), a projective morphism
\(\pi : X \to S\), and a coherent \(\OO_X\)-module \(\F\). The conclusion of Nitsure's
theorem covers an arbitrary noetherian base with a projective \(X \to S\), which
applies in particular to \(X = C \times_k T\), where \(C\) is a smooth proper
curve over a field \(k\), \(T\) is a noetherian \(k\)-scheme, and \(\F\) is a
coherent sheaf on the relative curve.

\section{Coherent flatness over a base}
\label{sec:coh_flat}

\begin{definition}[Coherent sheaf flat over the base]\leanok
  \label{def:coherent_sheaf_flat}
  \dcref{custom}
  \lean{AlgebraicGeometry.Scheme.CoherentSheafFlat}
  Let \(S\) be a noetherian scheme, \(f : X \to S\) a morphism of finite type, and
  \(\F\) a coherent \(\OO_X\)-module.  We say \(\F\) is \emph{flat over \(S\)}
  (equivalently, \(\F\) is \emph{\(f\)-flat}, or \emph{\(S\)-flat}) if for every
  point \(x \in X\) the stalk \(\F_x\) is a flat \(\OO_{S, f(x)}\)-module, where the
  \(\OO_{S, f(x)}\)-module structure is given by the composite
  \(\OO_{S, f(x)} \to \OO_{X, x} \to \End(\F_x)\).
  \textit{Source: cf.\ \cite{stacks-project}, tag 00HB (flat modules); \cite{nitsure-hilbert-quot}, \S4
  passim.}
  Equivalently (after pulling back to an affine open \(U = \Spec A \subseteq S\) and
  an affine open \(V = \Spec B \subseteq f^{-1}(U)\) with \(\F|_V = \widetilde{M}\)
  for \(M\) a finite \(B\)-module), \(\F|_V\) is \(\OO_U\)-flat in the sense that \(M\) is a
  flat \(A\)-module (with the \(A\)-module structure coming from the ring map
  \(A \to B\)).
\end{definition}

\begin{remark}
\label{rem:coh_flat_pullback}
  \dcref{custom}
  The condition is stable under base change: for any morphism \(g : S' \to S\),
  pulling back along the cartesian square
  \[
    \begin{tikzcd}
      X' \ar[r, "g'"] \ar[d, "f'"'] & X \ar[d, "f"] \\
      S' \ar[r, "g"]                & S
    \end{tikzcd}
  \]
  yields a coherent \(\OO_{X'}\)-module \(\F' := (g')^* \F\) which is \(S'\)-flat
  whenever \(\F\) is \(S\)-flat.  Indeed, at each stalk the pullback is an
  extension of scalars of the original flat module, and extension of scalars
  preserves flatness. The converse fails in general.
\end{remark}

\section{Generic flatness}
\label{sec:generic_flatness}

The construction of the flattening stratification proceeds by Noetherian
induction on \(S\). The induction step relies on the following classical
\emph{generic flatness} theorem of Grothendieck: over a noetherian integral
base, a coherent sheaf is flat over a non-empty open subscheme. Algebraically,
this is the assertion that a finitely generated module over a finite-type
algebra over a noetherian domain is generically free after inverting a single
non-zero element of the base ring.

\subsection{The dévissage layer}
\label{sec:generic_freeness_devissage}

The proof of the algebraic generic-freeness theorem has two layers: a
\emph{dévissage} reducing an arbitrary finite module over the (noetherian)
finite-type algebra \(B\) to the case of a quotient domain \(B/\pp\), and a
\emph{normalisation core} handling that case. We first develop the dévissage
layer as standalone module theory.

\begin{definition}[Generically free module]\leanok
  \label{def:generically_free}
  \dcref{custom}
  \lean{AlgebraicGeometry.GenericFreeness.GenericallyFree}
  Let \(A\) be a commutative ring and \(M\) an \(A\)-module.  We say \(M\) is
  \emph{generically free over \(A\)} if there exists \(f \in A\), \(f \ne 0\),
  such that the localisation \(M_f\) is a free \(A_f\)-module.
\end{definition}

\begin{lemma}[Transport of generic freeness along further localisation]\leanok
  \label{lem:generically_free_transport}
  \dcref{custom}
  \lean{AlgebraicGeometry.GenericFreeness.free_localizedModule_powers_mul}
  Let \(A\) be a commutative ring, \(M\) an \(A\)-module, and \(f, g \in A\).
  If \(M_f\) is a free \(A_f\)-module, then \(M_{fg}\) is a free
  \(A_{fg}\)-module.
\end{lemma}

\begin{proof}
\leanok
  Localisation is base change: \(M_h \cong A_h \otimes_A M\) as
  \(A_h\)-modules for any \(h \in A\).  Since \(fg\) maps to a unit multiple
  of \(g\) in \(A_{fg}\), the ring \(A_{fg}\) is canonically an
  \(A_f\)-algebra, and base-change associativity gives
  \[
    M_{fg} \;\cong\; A_{fg} \otimes_A M
    \;\cong\; A_{fg} \otimes_{A_f} \bigl(A_f \otimes_A M\bigr)
    \;\cong\; A_{fg} \otimes_{A_f} M_f .
  \]
  The base change of a free module along any ring map is free (a basis of
  \(M_f\) over \(A_f\) tensors to a basis over \(A_{fg}\)), so \(M_{fg}\) is
  free over \(A_{fg}\).
\end{proof}

\begin{lemma}[Generic freeness is invariant under isomorphism]\leanok
  \label{lem:generically_free_of_equiv}
  \dcref{custom}
  \lean{AlgebraicGeometry.GenericFreeness.GenericallyFree.of_linearEquiv}
  \uses{def:generically_free}
  Let \(A\) be a commutative ring and \(M \cong N\) isomorphic \(A\)-modules.
  If \(M\) is generically free over \(A\), so is \(N\).
\end{lemma}

\begin{proof}
\leanok
  An isomorphism \(e : M \to N\) localises to an isomorphism
  \(M_f \to N_f\) of \(A_f\)-modules for the witness \(f\), and freeness
  transports along it.
\end{proof}

\begin{lemma}[Splicing: generic freeness is stable under extensions]\leanok
  \label{lem:generically_free_splice}
  \lean{AlgebraicGeometry.GenericFreeness.genericallyFree_of_exact}
  \uses{def:generically_free}
  \dcref{custom}
  Let \(A\) be a domain and \(0 \to M_1 \to M_2 \to M_3 \to 0\) a short exact
  sequence of \(A\)-modules.  If \(M_1\) and \(M_3\) are generically free over
  \(A\), then so is \(M_2\).
\end{lemma}

\begin{proof}
\leanok
  \uses{lem:generically_free_transport}
  Let \(f_1, f_3 \ne 0\) be witnesses for \(M_1\), \(M_3\) and set
  \(f := f_1 f_3 \ne 0\) (here \(A\) is a domain).  Localisation at \(f\) is
  exact, so \(0 \to (M_1)_f \to (M_2)_f \to (M_3)_f \to 0\) is a short exact
  sequence of \(A_f\)-modules.  By
  \cref{lem:generically_free_transport}, \((M_1)_f\) and \((M_3)_f\) are free
  \(A_f\)-modules.  A short exact sequence with free (hence projective)
  quotient splits, so \((M_2)_f \cong (M_1)_f \oplus (M_3)_f\) is free.
\end{proof}

\begin{lemma}[Generic freeness of the zero module]\leanok
  \label{lem:generically_free_subsingleton}
  \dcref{custom}
  \lean{AlgebraicGeometry.GenericFreeness.genericallyFree_of_subsingleton}
  \uses{def:generically_free}
  Let \(A\) be a non-trivial commutative ring.  The zero \(A\)-module is
  generically free over \(A\), with witness \(f = 1\).
\end{lemma}

\begin{proof}
\leanok
  The localisation of the zero module is zero, which is free (on the empty
  basis), and \(1 \ne 0\) since \(A\) is non-trivial.
\end{proof}

\begin{lemma}[Torsion base case]\leanok
  \label{lem:generically_free_torsion}
  \lean{AlgebraicGeometry.GenericFreeness.genericallyFree_of_torsion}
  \uses{def:generically_free}
  \dcref{custom}
  Let \(A\) be a domain and \(M\) a finite \(A\)-module every element of which
  is annihilated by some non-zero element of \(A\).  Then \(M\) is generically
  free over \(A\): indeed some localisation \(M_f\) with \(f \ne 0\) vanishes.
\end{lemma}

\begin{proof}
\leanok
  Choose a finite generating set \(m_1, \dots, m_r\) of \(M\) and non-zero
  annihilators \(a_i m_i = 0\).  The product \(f := a_1 \cdots a_r\) is
  non-zero (\(A\) is a domain) and annihilates every generator (extract the
  factor \(a_i\) and commute), hence — the annihilator of \(f\) being a
  submodule containing the generators — annihilates all of \(M\).  Every
  element of \(M_f\) is then zero, and the zero module is free.
\end{proof}

\begin{lemma}[Generic freeness: finite modules]\leanok
  \label{lem:generically_free_finite}
  \dcref{custom}
  \lean{AlgebraicGeometry.GenericFreeness.exists_free_localizationAway_of_finite}
  \uses{def:generically_free}
  Let \(A\) be a noetherian domain and \(M\) a finite \(A\)-module.  Then \(M\)
  is generically free over \(A\).
\end{lemma}

\begin{proof}
\leanok
  A finite module over a noetherian ring is finitely presented, and
  \(M \otimes_A \mathrm{Frac}(A)\) is a finite-dimensional
  \(\mathrm{Frac}(A)\)-vector space, hence free.  Freeness at the generic
  point of a finitely presented module spreads out to a basic open: there is
  \(f \ne 0\) in \(A\) with \(M_f\) free over \(A_f\).
\end{proof}

\begin{lemma}[Generic flatness: finite modules]\leanok
  \label{lem:generically_flat_finite}
  \dcref{custom}
  \lean{AlgebraicGeometry.GenericFreeness.exists_flat_localizationAway_of_finite}
  \uses{def:generically_free}
  Let \(A\) be a noetherian domain and \(M\) a finite \(A\)-module.  Then there
  is \(f \ne 0\) in \(A\) with \(M_f\) flat over \(A_f\).
\end{lemma}

\begin{proof}
\leanok
  \uses{lem:generically_free_finite}
  Immediate from \cref{lem:generically_free_finite}: free modules are flat.
\end{proof}

\begin{lemma}[Generic freeness: the module-finite case]\leanok
  \label{lem:generically_free_module_finite}
  \dcref{custom}
  \lean{AlgebraicGeometry.GenericFreeness.exists_free_localizationAway_of_moduleFinite}
  \uses{def:generically_free}
  Let \(A\) be a noetherian domain, \(B\) a module-finite \(A\)-algebra, and
  \(M\) a finite \(B\)-module.  Then \(M\) is generically free over \(A\).
  (This is the case of \cref{thm:generic_flatness_algebraic} where the
  morphism \(\Spec B \to \Spec A\) is finite.)
\end{lemma}

\begin{proof}
\leanok
  \uses{lem:generically_free_finite}
  A finite module over a module-finite \(A\)-algebra is a finite
  \(A\)-module, so \cref{lem:generically_free_finite} applies.
\end{proof}

\begin{lemma}[Generic-rank comparison for a finite module over a domain]\leanok
  \label{lem:generic_rank_comparison}
  \dcref{custom}
  \lean{AlgebraicGeometry.GenericFreeness.exists_generic_rank_comparison}
  Let \(D\) be a domain and \(M\) a finite \(D\)-module.  Then there exist
  \(r \in \mathbb{N}\), an injective \(D\)-linear map
  \(\Phi : D^{\oplus r} \to M\),
  and a non-zero \(d \in D\) annihilating the cokernel
  \(M / \Phi(D^{\oplus r})\); that is, \(M\) sits in an exact sequence
  \(0 \to D^{\oplus r} \to M \to T \to 0\) with \(d\,T = 0\).
\end{lemma}

\begin{proof}
\leanok
  Localise at \(S := D \setminus \{0\}\): the localisation \(S^{-1}M\) is a
  finite-dimensional vector space over \(\mathrm{Frac}(D) = S^{-1}D\); let
  \(r\) be its dimension.  Every element of \(S^{-1}M\) has the form
  \(m/s\), so scaling each vector of a basis by the (unit) image of its
  denominator produces a basis \((\bar c_1, \dots, \bar c_r)\) whose members
  are images of elements \(c_1, \dots, c_r \in M\); let
  \(\Phi(e_i) := c_i\).  A \(D\)-linear relation among the \(c_i\)
  localises to a \(\mathrm{Frac}(D)\)-linear relation among the basis
  vectors \(\bar c_i\), whose coefficients — images of the original ones —
  must vanish, and \(D \to \mathrm{Frac}(D)\) is injective, so \(\Phi\) is
  injective.  For the cokernel: given \(y \in M\), expand its image in the
  basis with coefficients in \(\mathrm{Frac}(D)\) and clear their
  denominators by a single \(s \in S\); then \(s \cdot y - \Phi(u)\) is in
  the kernel of \(M \to S^{-1}M\) for some \(u \in D^{\oplus r}\), hence is
  killed by some \(t \in S\), so \(ts \ne 0\) annihilates the class of
  \(y\) in the cokernel \(T\).  Finally \(T\) is a finite \(D\)-module, so
  the product of such annihilators over a finite generating set is a single
  non-zero uniform annihilator \(d\).
\end{proof}

\begin{lemma}[Normalisation core: generic freeness of quotient domains]\leanok
  \label{lem:generically_free_domain_core}
  \lean{AlgebraicGeometry.GenericFreeness.genericallyFree_quotient_prime}
  \uses{def:generically_free}
  \dcref{custom}
  Let \(A\) be a noetherian domain, \(B\) a finite-type \(A\)-algebra, and
  \(\pp \subset B\) a prime ideal.  Then the quotient domain \(B/\pp\) is
  generically free as an \(A\)-module.
\end{lemma}

\begin{proof}
\leanok
  \uses{lem:generically_free_torsion, lem:generically_free_transport,
    lem:generically_free_splice, lem:generic_rank_comparison}
  Write \(C := B/\pp\), a finite-type \(A\)-algebra which is a domain, and let
  \(K := \mathrm{Frac}(A)\).  If the kernel of \(A \to C\) is non-zero, every
  element of \(C\) is annihilated by any non-zero \(a \in \ker(A \to C)\)
  (as \(a \cdot c = (a \cdot 1)\, c = 0\)), and
  \cref{lem:generically_free_torsion} concludes.  So assume \(A \subseteq C\).
  Then \(C_K := K \otimes_A C\) is a non-zero finite-type \(K\)-algebra domain;
  let \(m\) be its Krull dimension.  We induct on \(m\), the statement for all
  smaller values (over all noetherian domains \(A\)) being granted.

  By Noether normalisation there are \(b_1, \dots, b_m \in C_K\), algebraically
  independent over \(K\), with \(C_K\) module-finite over
  \(K[b_1, \dots, b_m]\); scaling by non-zero elements of \(A\) (which does not
  affect either property) we may take \(b_i \in C\).  Fix a finite generating
  set of \(C\) as an \(A\)-algebra; each generator satisfies a monic equation
  over \(K[b_1, \dots, b_m]\), and clearing the finitely many denominators of
  all coefficients produces a single \(g \ne 0\) in \(A\) such that each
  generator of \(C_g\) is integral over \(A_g[b_1, \dots, b_m]\) (multiply the
  monic equation for \(x\) through by \(g^{\deg}\) and rewrite it as a monic
  equation for a unit multiple of \(x\)).  A finite-type algebra which is
  integral is module-finite, so \(C_g\) is a finite module over the polynomial
  ring \(D := A_g[b_1, \dots, b_m]\).

  Let \(r\) be the dimension of \(C_g \otimes_D \mathrm{Frac}(D)\) as a
  \(\mathrm{Frac}(D)\)-vector space and choose \(c_1, \dots, c_r \in C_g\)
  mapping to a basis.  The resulting map \(D^{\oplus r} \to C_g\) is injective
  (a \(D\)-linear relation among the \(c_i\) becomes a
  \(\mathrm{Frac}(D)\)-linear relation), and its cokernel \(T\) is a finite
  \(D\)-module with \(T \otimes_D \mathrm{Frac}(D) = 0\), i.e.\ torsion;
  this generic-rank packaging, including the uniform annihilator below, is
  \cref{lem:generic_rank_comparison}.  A
  non-zero \(d \in D\) annihilating \(T\) exhibits
  \(T_K := K \otimes_{A_g} T\) as a finite module over
  \(K[b_1, \dots, b_m]/(d')\) for the non-zero image \(d'\) of \(d\), a ring of
  Krull dimension \(< m\); presenting that ring as a quotient of a polynomial
  ring and filtering as in \cref{thm:generic_flatness_algebraic}'s dévissage,
  the inductive hypothesis (over the noetherian domain \(A_g\)) together with
  \cref{lem:generically_free_splice} yields \(h \ne 0\) in \(A_g\), which we
  may take in \(A\), with \(T_h\) free over \((A_g)_h\).  Localising the exact
  sequence \(0 \to D^{\oplus r} \to C_g \to T \to 0\) at \(h\) and splitting
  off the free (hence projective) quotient \(T_h\), we get
  \[
    C_{gh} \;\cong\; \bigl(A_{gh}[b_1, \dots, b_m]\bigr)^{\oplus r} \oplus T_h ,
  \]
  a direct sum of free \(A_{gh}\)-modules (a polynomial ring is free on its
  monomials), so \(f := gh\) is a generic-freeness witness for \(C\)
  (\cref{lem:generically_free_transport} normalising the witnesses).
\end{proof}

\begin{theorem}[Generic flatness, algebraic form]\leanok
  \label{thm:generic_flatness_algebraic}
  \lean{AlgebraicGeometry.GenericFreeness.genericFlatnessAlgebraic}
  \uses{def:generically_free}
  \dcref{custom}
  \textit{Source: \cite{nitsure-hilbert-quot}, \S4, ``Lemma on Generic Flatness''.}
  Let \(A\) be a noetherian domain, \(B\) a finite-type \(A\)-algebra, and \(M\) a finite
  \(B\)-module. Then \(M\) is generically free over \(A\): there exists
  \(f \in A\), \(f \ne 0\), such that the localisation
  \(M_f\) is a free \(A_f\)-module.
\end{theorem}

\begin{proof}
\leanok
  \uses{lem:generically_free_subsingleton, lem:generically_free_of_equiv,
    lem:generically_free_splice, lem:generically_free_domain_core}
  The ring \(B\) is noetherian (it is of finite type over the noetherian
  \(A\)), so the finite \(B\)-module \(M\) admits a finite prime filtration
  \[
    0 = M_0 \subset M_1 \subset \cdots \subset M_r = M
  \]
  with successive quotients \(M_i / M_{i-1} \cong B / \mathfrak{p}_i\) for
  prime ideals \(\mathfrak{p}_i \subset B\).  By induction on the filtration
  length it suffices that generic freeness over \(A\) holds for the zero
  module (\cref{lem:generically_free_subsingleton}), holds for each quotient
  domain \(B/\mathfrak{p}\) (\cref{lem:generically_free_domain_core},
  transported along the isomorphism \(M_i/M_{i-1} \cong B/\mathfrak{p}_i\) by
  \cref{lem:generically_free_of_equiv}), and is stable under extensions
  (the splicing lemma, \cref{lem:generically_free_splice}, applied to
  \(0 \to M_{i-1} \to M_i \to M_i/M_{i-1} \to 0\)).
\end{proof}

\section{Geometric glue: from generic freeness to generic flatness}

The algebraic statement concerns one module over one algebra; the geometric
statement quantifies over all affine pairs below the witness open. The glue
consists of a mixed-base stability lemma for flatness under module
localization, two restriction-compatibility lemmas for the induced module
structures, and a two-layer basic-open reduction.

\begin{lemma}[Mixed-base stability of flatness under module localization]
  \leanok
  \label{lem:flat_localizedModule_algebra}
  \dcref{custom}
  \lean{Module.Flat.of_isLocalizedModule_algebra}
  Let \(B\) be an algebra over a commutative ring \(R\), let \(N\) be a
  \(B\)-module which is flat as an \(R\)-module, and let
  \(g : N \to N'\) be a localization of \(B\)-modules at a submonoid
  \(p \subseteq B\). Then \(N'\) is flat as an \(R\)-module.
\end{lemma}
\begin{proof}
  \leanok
  Let \(f : Q \hookrightarrow Q'\) be an injective map of \(R\)-modules. The
  square
  \[
    \begin{array}{ccc}
      N \otimes_R Q & \longrightarrow & N \otimes_R Q' \\
      \downarrow & & \downarrow \\
      N' \otimes_R Q & \longrightarrow & N' \otimes_R Q'
    \end{array}
  \]
  identifies the bottom map \(\mathrm{id}_{N'} \otimes f\) with the
  localization at \(p\) of the top map \(\mathrm{id}_N \otimes f\): indeed
  \(N' \otimes_R Q\) is the localization of the \(B\)-module
  \(N \otimes_R Q\) at \(p\) (localization commutes with \(\, \otimes_R Q\),
  because \(S^{-1}(N \otimes_R Q) = S^{-1}B \otimes_B (N \otimes_R Q)
  = (S^{-1}B \otimes_B N) \otimes_R Q\)), and localizing the top map yields
  the bottom one. The top map is injective by flatness of \(N\) over \(R\),
  and localization of modules preserves injectivity; hence
  \(\mathrm{id}_{N'} \otimes f\) is injective.
\end{proof}

\begin{lemma}[Fibre-side restriction compatibility of \(\mathrm{appLE}\)]
  \leanok
  \label{lem:appLE_res_apply}
  \dcref{custom}
  \lean{AlgebraicGeometry.appLE_res_apply}
  Let \(p : X \to S\) be a morphism of schemes, \(U \subseteq S\) open,
  \(W' \subseteq W \subseteq p^{-1}(U)\) opens of \(X\), and
  \(u \in \Gamma(S, U)\). Then restricting
  \(p^\sharp_{U,W}(u) \in \Gamma(X, W)\) to \(W'\) yields
  \(p^\sharp_{U,W'}(u)\).
\end{lemma}
\begin{proof}
  \leanok
  Both sides are the composite of \(p^\sharp\) with restriction along
  \(W' \subseteq W \subseteq p^{-1}(U)\); this is the naturality of
  \(p^\sharp\) in the target open.
\end{proof}

\begin{lemma}[Base-side restriction compatibility of \(\mathrm{appLE}\)]
  \leanok
  \label{lem:appLE_base_res_apply}
  \dcref{custom}
  \lean{AlgebraicGeometry.appLE_base_res_apply}
  Let \(p : X \to S\) be a morphism of schemes, \(U' \subseteq U \subseteq S\)
  opens, \(W \subseteq p^{-1}(U')\) an open of \(X\), and
  \(u \in \Gamma(S, U)\). Then
  \(p^\sharp_{U',W}\bigl(u|_{U'}\bigr) = p^\sharp_{U,W}(u)\).
\end{lemma}
\begin{proof}
  \leanok
  Naturality of \(p^\sharp\) in the source open.
\end{proof}

\begin{lemma}[Section flatness transports along equal opens]
  \leanok
  \label{lem:flat_sections_congr}
  \dcref{custom}
  \lean{AlgebraicGeometry.flat_sections_congr}
  Let \(p : X \to S\), \(\mathcal F\) a module sheaf on \(X\),
  \(U_0 \subseteq S\) open, and \(O_1 = O_2 \subseteq p^{-1}(U_0)\) equal
  opens of \(X\). If \(\Gamma(\mathcal F, O_1)\) is flat over
  \(\Gamma(S, U_0)\) (via \(p^\sharp\)), then so is
  \(\Gamma(\mathcal F, O_2)\).
\end{lemma}
\begin{proof}
  \leanok
  Immediate: the two data coincide.
\end{proof}

\begin{lemma}[Chart-basic core of generic flatness]
  \leanok
  \label{lem:flat_section_chartBasic}
  \dcref{custom}
  \lean{AlgebraicGeometry.flat_section_chartBasic}
  \uses{def:coherent_sheaf_flat}
  Let \(p : X \to S\), let \(\mathcal F\) be quasi-coherent on \(X\), let
  \(U_0 \subseteq S\) be open, and let \(W_j \subseteq p^{-1}(U_0)\) be an
  affine chart. Let \(g \in A := \Gamma(S, U_0)\) be such that the localized
  section module \(\Gamma(\mathcal F, W_j)_g\) is free over \(A_g\). Then for
  every \(c \in \Gamma(X, W_j)\) with \(D(c) \subseteq p^{-1}(D(g))\), the
  module \(\Gamma(\mathcal F, D(c))\) is flat over \(A\).
\end{lemma}
\begin{proof}
  \leanok
  \uses{lem:flat_localizedModule_algebra, lem:appLE_res_apply,
    lem:flat_sections_congr, lem:qcoh_section_localization_basicOpen}
  Write \(\beta := p^\sharp_{U_0, W_j}(g) \in B_j := \Gamma(X, W_j)\), so that
  \(D(\beta) = W_j \cap p^{-1}(D(g))\). By the quasi-compact--quasi-separated
  section-localization theorem
  (\cref{lem:qcoh_section_localization_basicOpen}, applied to the affine
  \(W_j\)), the restriction
  \(\Gamma(\mathcal F, W_j) \to \Gamma(\mathcal F, D(\beta))\) is the
  localization of \(B_j\)-modules at the powers of \(\beta\); since
  \(\beta^k\) acts on these modules as \(g^k\) acts through \(p^\sharp\), it
  is also the localization of \(A\)-modules at the powers of \(g\). Hence
  \(\Gamma(\mathcal F, D(\beta)) \cong \Gamma(\mathcal F, W_j)_g\) is free,
  in particular flat, over \(A_g\), and therefore flat over \(A\) (flatness
  over \(A\) and over the localization \(A_g\) agree for
  \(A_g\)-modules).

  Since \(D(c) \subseteq D(\beta)\), the further restriction
  \(\Gamma(\mathcal F, D(\beta)) \to \Gamma(\mathcal F, D(c))\) is again a
  section localization (\cref{lem:qcoh_section_localization_basicOpen} on the
  affine \(D(\beta)\), at the restricted element \(c|_{D(\beta)}\), whose
  basic open is \(D(c)\)) — a localization of
  \(\Gamma(X, D(\beta))\)-modules. Mixed-base stability
  (\cref{lem:flat_localizedModule_algebra}, with base \(A\) and fibre algebra
  \(\Gamma(X, D(\beta))\)) shows \(\Gamma(\mathcal F, D(c|_{D(\beta)}))\) is
  flat over \(A\); the induced-structure compatibilities are
  \cref{lem:appLE_res_apply}, and \cref{lem:flat_sections_congr} transports
  the conclusion along \(D(c|_{D(\beta)}) = D(c)\).
\end{proof}

\begin{lemma}[Two-layer basic-open reduction]
  \leanok
  \label{lem:flat_section_pair}
  \dcref{custom}
  \lean{AlgebraicGeometry.flat_section_pair}
  \uses{def:coherent_sheaf_flat}
  Let \(p : X \to S\), \(\mathcal F\) quasi-coherent on \(X\),
  \(U_0 \subseteq S\) an affine open with \(A := \Gamma(S, U_0)\), and let
  \(\{W_c^{(j)}\}_{j}\) be affine charts contained in and covering
  \(p^{-1}(U_0)\). Let \(f \in A\) be such that each localized section
  module \(\Gamma(\mathcal F, W_c^{(j)})_f\) is free over \(A_f\). Then for
  every affine pair \(U \subseteq D(f)\), \(W \subseteq p^{-1}(U)\), the
  module \(\Gamma(\mathcal F, W)\) is flat over \(\Gamma(S, U)\) via
  \(p^\sharp\).
\end{lemma}
\begin{proof}
  \leanok
  \uses{lem:flat_section_chartBasic, lem:appLE_base_res_apply,
    lem:qcoh_section_localization_basicOpen, lem:generically_free_transport}
  For each \(x \in W\): choose (affine interchange on \(S\), applied to the
  affines \(U\) and \(U_0\) at the point \(p(x) \in U \cap U_0\)) sections
  \(a \in \Gamma(S, U)\), \(a' \in A\) with
  \(D(a) = D(a') \ni p(x)\); choose a chart \(W_c^{(j)} \ni x\); shrink to a
  basic open of \(W\) inside \(W \cap p^{-1}(D(a)) \cap W_c^{(j)}\)
  containing \(x\); and apply affine interchange on \(X\) to that basic open
  and the affine \(W_c^{(j)}\) to obtain \(b \in \Gamma(X, W)\),
  \(c \in \Gamma(X, W_c^{(j)})\) with
  \(x \in D(b) = D(c) \subseteq p^{-1}(D(a))\).

  The basic opens \(D(b)\) obtained this way cover \(W\), so the chosen
  elements \(b\) generate the unit ideal of \(\Gamma(X, W)\). Since each
  restriction \(\Gamma(\mathcal F, W) \to \Gamma(\mathcal F, D(b))\) is a
  localization of \(\Gamma(X, W)\)-modules at the powers of \(b\)
  (\cref{lem:qcoh_section_localization_basicOpen}), flatness of
  \(\Gamma(\mathcal F, W)\) over \(\Gamma(S, U)\) may be checked on the
  pieces \(\Gamma(\mathcal F, D(b))\) (flatness over the base is local with
  respect to a spanning family of the fibre algebra).

  Fix a piece, with its data \(a, a', j, c\). Freeness of
  \(\Gamma(\mathcal F, W_c^{(j)})_f\) localizes further to
  \(f a'\) (\cref{lem:generically_free_transport}), and
  \(D(c) = D(b) \subseteq p^{-1}(D(a)) = p^{-1}(D(a'))
  \subseteq p^{-1}(D(f a'))\) since \(D(a) \subseteq U \subseteq D(f)\).
  The chart-basic core (\cref{lem:flat_section_chartBasic}, at the element
  \(f a'\) and the chart section \(c\)) gives flatness of
  \(\Gamma(\mathcal F, D(b))\) over \(A\). Finally exchange the base ring
  twice along the localizations \(A \to \Gamma(S, D(a')) = \Gamma(S, D(a))\)
  and \(\Gamma(S, U) \to \Gamma(S, D(a))\) (for a module over a localized
  ring, flatness over the ring and over its localization agree; the
  compatibility of the induced structures is
  \cref{lem:appLE_base_res_apply}), obtaining flatness of
  \(\Gamma(\mathcal F, D(b))\) over \(\Gamma(S, U)\).
\end{proof}

\begin{theorem}[Generic flatness, geometric form]
  \leanok
  \label{thm:generic_flatness}
  \dcref{custom}
  \lean{AlgebraicGeometry.genericFlatness}
  \uses{thm:generic_flatness_algebraic}
  \textit{Source: \cite{nitsure-hilbert-quot}, \S4, Theorem ``generic flatness''.}
  Let \(S\) be an integral locally noetherian scheme, \(p : X \to S\) a
  finite-type morphism (that is, quasi-compact and locally of finite type —
  quasi-compactness is essential: it bounds the number of denominators to
  clear, and the statement fails for the locally-of-finite-type morphism
  \(\coprod_q \Spec \mathbb{F}_q \to \Spec \Z\)), and \(\F\) a coherent
  \(\OO_X\)-module. Then there exists a non-empty open subscheme
  \(U \subseteq S\) such that the restriction of \(\F\) to
  \(X_U = p^{-1}(U)\) is \(\OO_U\)-flat.
\end{theorem}

\begin{proof}
  \leanok
  \uses{thm:generic_flatness_algebraic, lem:flat_section_pair,
    lem:affine_finite_sections_nhds}
  Choose a non-empty affine open \(U_0 \subseteq S\); its section ring
  \(A := \Gamma(S, U_0)\) is a noetherian domain because \(S\) is integral and
  locally noetherian. The preimage \(p^{-1}(U_0)\) is quasi-compact
  (\(p\) is quasi-compact and \(U_0\) is), and every point of it has an
  affine open neighbourhood inside \(p^{-1}(U_0)\) on which the sections of
  \(\F\) are a finite module over the sections of \(\OO_X\)
  (\cref{lem:affine_finite_sections_nhds}); extract a finite subcover by
  affine charts \(W^{(1)}, \dots, W^{(r)}\). Each
  \(B_i := \Gamma(X, W^{(i)})\) is of finite type over \(A\) (as \(p\) is
  locally of finite type and both opens are affine) and each
  \(M_i := \Gamma(\F, W^{(i)})\) is a finite \(B_i\)-module, so algebraic
  generic freeness (\cref{thm:generic_flatness_algebraic}) produces
  \(f_i \in A \setminus \{0\}\) with \((M_i)_{f_i}\) free over \(A_{f_i}\).
  Set \(f := f_1 \cdots f_r \ne 0\); freeness localizes from each \(f_i\) to
  \(f\), so every \((M_i)_f\) is free over \(A_f\). The witness open is
  \(V := D(f) \subseteq U_0\), non-empty because \(S\) is reduced and
  \(f \ne 0\). The required flatness on every affine pair
  \(U \subseteq V\), \(W \subseteq p^{-1}(U)\) is exactly the two-layer
  basic-open reduction, \cref{lem:flat_section_pair}.
\end{proof}

\section{Existence of the flattening stratification}
\label{sec:flatstrat_exists}

\begin{theorem}

[Existence of the flattening stratification]
  \label{thm:flattening_stratification_exists}
  \dcref{custom}
  \uses{def:coherent_sheaf_flat, lem:flat_locus_open, lem:nonflat_locus_proper,
    lem:noetherian_induction_strata, thm:generic_flatness}
  \textit{Source: \cite{nitsure-hilbert-quot}, \S4, Theorem ``flattening stratification'';
  cf.\ \cite{stacks-project}, tag 052H.}
  Let \(S\) be a noetherian scheme and \(\F\) a coherent sheaf on the relative
  projective space \(\PP^n_S\) over \(S\).  Then:
  \begin{enumerate}[label=(\roman*)]
    \item The set \(I\) of Hilbert polynomials
      \(P_s(m) = \chi(\PP^n_s, \F_s(m))\) that arise as \(s\) ranges over \(S\) is
      finite.
    \item For each \(f \in I\) there exists a locally-closed subscheme
      \(S_f \hookrightarrow S\) whose underlying set \(|S_f|\) consists of exactly
      those \(s \in S\) with \(P_s = f\).  The collection \(\{|S_f|\}_{f \in I}\) is
      pairwise disjoint with set-theoretic union \(|S|\).  Forming the coproduct
      \(S' := \coprod_{f \in I} S_f\) and the morphism \(i : S' \to S\) assembled
      from the inclusions \(S_f \hookrightarrow S\), the pullback \(i^* \F\) on
      \(\PP^n_{S'}\) is flat over \(S'\), and \(i : S' \to S\) is universal with this
      property: for any morphism \(\varphi : T \to S\) such that \(\varphi^* \F\) on
      \(\PP^n_T\) is flat over \(T\), \(\varphi\) factors uniquely through \(i\).
      Each \(S_f\) is uniquely determined by \(f\).
    \item Order \(I\) by \(f < g \iff f(n) < g(n)\) for all \(n \gg 0\). Then the
      closure of \(|S_f|\) in \(|S|\) is contained in \(\bigcup_{g \ge f} |S_g|\).
  \end{enumerate}
\end{theorem}

\begin{proof}
  \uses{def:coherent_sheaf_flat, lem:flat_locus_open, lem:nonflat_locus_proper,
    lem:noetherian_induction_strata, thm:generic_flatness}
  Following \cite{nitsure-hilbert-quot}, \S4, reduce to the special case \(n = 0\) and assemble
  the resulting strata over the relative twists \(\pi_* \F(N+i)\).  The
  construction is local on \(S\), since the strata glue by their universal
  property.  We establish the three conclusions in order.

  \emph{Part (i): finiteness of the Hilbert-polynomial set.}
  By Lemma~\ref{lem:nonflat_locus_proper}, Noetherian induction on \(S\) (whose
  inductive step invokes generic flatness, Theorem~\ref{thm:generic_flatness},
  on the integral locally-closed strata cut out by the irreducible-component
  decomposition of \(S\)) produces finitely many reduced, pairwise disjoint,
  locally-closed subschemes \(V_i \hookrightarrow S\) set-theoretically covering
  \(|S|\) on each of which \(\F|_{\PP^n_{V_i}}\) is flat over \(\OO_{V_i}\) in
  the sense of Definition~\ref{def:coherent_sheaf_flat}.  On a flat family the
  Hilbert polynomial \(P_s(m) = \chi(\PP^n_s, \F_s(m))\) is locally constant, so
  it takes finitely many values on each quasi-compact \(V_i\); the finite union
  over the finitely many \(V_i\) shows that the set \(I\) of Hilbert polynomials
  occurring on \(S\) is finite.

  \emph{Part (ii): existence of the strata and their universal property.}
  Working over the flat pieces \(V_i\) of part (i), semicontinuity
  (cohomology-and-base-change) yields a uniform integer \(N\) such that, for all
  \(s \in S\) and \(m \ge N\), the higher cohomologies \(H^r(\PP^n_s, \F_s(m))\)
  vanish for \(r \ge 1\) and the base-change homomorphism
  \((\pi_* \F(m))|_s \to H^0(\PP^n_s, \F_s(m))\) is an isomorphism.  With this
  \(N\) fixed, Lemma~\ref{lem:noetherian_induction_strata} applies the \(n = 0\)
  stratification mechanism (Lemma~\ref{lem:flat_locus_open}) successively to the
  coherent sheaves \(E_i = \pi_* \F(N+i)\) on \(S\) for \(i = 0, \dots, n\),
  refining stratum by stratum.  After \(n+1\) steps this produces a finite
  family of locally-closed subschemes \(W_{e_0, \dots, e_n} \subseteq S\),
  universal for the property that each \(E_i\) pulls back to a locally-free sheaf
  of constant rank \(e_i\); re-indexing the tuple \((e_0, \dots, e_n)\) by the
  unique numerical polynomial \(f\) of degree \(\le n\) with \(f(N+i) = e_i\)
  gives the family \(\{W_f\}_{f \in I}\).  At any \(s \in W_f\) the base-change
  isomorphism and the cohomology vanishing force \(P_s = f\) (since \(P_s\) has
  degree \(\le n\) and is determined by its values at \(N, \dots, N+n\)), so
  \(|W_f| = \{s : P_s = f\}\); these subsets are disjoint with union \(|S|\).
  The required locally-closed subscheme \(S_f \subseteq W_f\) is the closed
  subscheme cut out by the further local-freeness conditions
  ``\(\pi_* \F(N+i)|_{W_f}\) is locally free of rank \(f(N+i)\) for all
  \(i \ge 0\)'' (a chain of coherent ideal sheaves that stabilises by
  Noetherianity), and \(\F|_{\PP^n_{S_f}}\) is flat over \(S_f\)
  (Definition~\ref{def:coherent_sheaf_flat}) because the graded module
  \(\bigoplus_{i \gg 0} \pi_* \F(i)|_{S_f}\) is locally free.  Forming
  \(S' = \coprod_{f \in I} S_f\) and the assembled morphism \(i : S' \to S\),
  the pullback \(i^* \F\) is \(S'\)-flat; and any \(\varphi : T \to S\) with
  \(\varphi^* \F\) flat over \(T\) has locally constant Hilbert polynomial, hence
  decomposes \(T = \coprod_f T_f\) into the open-closed loci where the polynomial
  is \(f\), each factoring uniquely through the corresponding \(S_f\) by the
  local-freeness universal property of Lemma~\ref{lem:flat_locus_open}.  This
  gives the unique factorisation through \(i\), and shows each \(S_f\) is
  uniquely determined by \(f\).

  \emph{Part (iii): closure of strata.}
  Order \(I\) by \(f < g \iff f(n) < g(n)\) for all \(n \gg 0\).  Since each
  \(S_f\) is a stratum of the \(n = 0\) flattening stratification applied to a
  single direct image \(\pi_* \F(p)\) for \(p \gg 0\) separating all polynomials
  in \(I\), the closure-of-strata assertion reduces to the upper-semicontinuity
  statement already established in the \(n = 0\) case
  (Lemma~\ref{lem:flat_locus_open}): the closure of \(|S_f|\) in \(|S|\) is
  contained in \(\bigcup_{g \ge f} |S_g|\).
\end{proof}

The construction begins with the special case \(n = 0\), where
\(\PP^n_S = S\) and \(\F\) is a coherent sheaf on \(S\).  Its remaining
ingredients are the Noetherian-induction reduction to the flat case, uniform
vanishing of the higher cohomology of \(\F(m)\) for \(m \gg 0\) on the flat
pieces, and base-change isomorphisms for the direct images \(\pi_*\F(m)\).

\begin{lemma}
[Flat locus on \(S\) for a coherent sheaf is a locally-closed stratification]
  \label{lem:flat_locus_open}
  \dcref{custom}
  \textit{Source: \cite{nitsure-hilbert-quot}, \S4 (special case \(n=0\) of the main theorem).}
  Let \(S\) be a noetherian scheme and let \(\F\) be a coherent \(\OO_S\)-module.
  For each integer \(e \ge 0\), there exists a locally-closed subscheme
  \(S_e \hookrightarrow S\) whose underlying set is
  \(\{s \in S : \dim_{\kappa(s)} \F|_s = e\}\). The family \(\{S_e\}_e\) is a
  finite, mutually disjoint, set-theoretically covering stratification of \(S\),
  and a morphism \(f : T \to S\) pulls back \(\F\) to a locally-free
  \(\OO_T\)-module of rank \(e\) if and only if \(f\) factors through \(S_e\).  The
  function \(s \mapsto \dim_{\kappa(s)} \F|_s\) is upper-semicontinuous on \(S\).
\end{lemma}

\begin{proof}
  By \cite{nitsure-hilbert-quot}, \S4, for \(s \in S\), Nakayama produces a
  neighbourhood \(V\) of \(s\) and a right-exact local presentation
  \(\OO_V^{\oplus m} \xrightarrow{\psi} \OO_V^{\oplus e} \to \F|_V \to 0\) where
  \(e = \dim_{\kappa(s)} \F|_s\).  The closed subscheme \(V_e \hookrightarrow V\)
  defined by the ideal sheaf generated by the matrix entries of \(\psi\)
  represents the locus where \(\F\) becomes locally free of rank \(e\): by
  right-exactness of tensor product, a base change \(f : T \to V\) has
  \(f^* \F\) locally free of rank \(e\) iff \(f^* \psi = 0\), iff \(f\) factors
  through \(V_e\).  Glueing the \(V_e\) over the cover of \(S\) by such
  neighbourhoods (and intersecting with \(V \setminus \bigcup_{e' > e} V_{e'}\)
  to get the strictly-equal-to-\(e\) stratum) produces the locally-closed
  \(S_e\). Upper-semicontinuity of \(e\) follows because the rank of \(\psi\) is
  lower-semicontinuous.
\end{proof}

\begin{lemma}
[Noetherian induction reduction to the flat case]
  \label{lem:nonflat_locus_proper}
  \dcref{custom}
  \uses{thm:generic_flatness}
  \textit{Source: \cite{nitsure-hilbert-quot}, \S4 (general case, reduction by Noetherian induction).}
  Let \(S\) be a noetherian scheme, \(\pi : \PP^n_S \to S\) the relative projective
  space, and \(\F\) a coherent sheaf on \(\PP^n_S\).  Then there exist finitely many
  reduced, locally-closed, pairwise disjoint subschemes \(V_i \hookrightarrow S\)
  whose set-theoretic union is \(|S|\) and such that \(\F|_{\PP^n_{V_i}}\) is
  flat over \(\OO_{V_i}\) for each \(i\). Consequently, only finitely many
  Hilbert polynomials \(P_s(m) = \chi(\PP^n_s, \F_s(m))\) occur as \(s\) varies
  over \(S\).
\end{lemma}

\begin{proof}
  By \cite{nitsure-hilbert-quot}, \S4: \(S\) being noetherian, decompose \(|S| = \bigcup Y_j\) into
  finitely many irreducible components, each closed.  For an irreducible
  component \(Y\), let \(U \subseteq Y\) be the open subset not lying on any other
  component, given the reduced subscheme structure (an integral, locally-closed
  subscheme of \(S\)).  By \cref{thm:generic_flatness} applied to
  \(\F|_{\PP^n_U}\), there is a non-empty open \(V \subseteq U\) over which \(\F\) is
  flat.  Replace \(S\) by the reduced closed complement \(S \setminus V\) and
  repeat. Noetherianity of \(S\) terminates the induction in finitely many
  steps, producing the \(V_i\).  On each \(V_i\) the Hilbert polynomial is locally
  constant (flat families have locally constant Hilbert polynomial), hence
  takes finitely many values on the noetherian (hence quasi-compact) scheme
  \(V_i\); the finite union of finite sets is finite.
\end{proof}

\begin{lemma}
[Assembly of the strata by Noetherian induction]
  \label{lem:noetherian_induction_strata}
  \dcref{custom}
  \uses{lem:flat_locus_open, lem:nonflat_locus_proper, thm:flat_base_change_cohomology}
  \textit{Source: \cite{nitsure-hilbert-quot}, \S4 (assembly step combining the \(n=0\) special case
  applied to direct images \(\pi_*\F(N+i)\)).}
  Let \(S\), \(\pi : \PP^n_S \to S\), \(\F\) be as in
  \cref{thm:flattening_stratification_exists}.  There exists an integer
  \(N \ge 0\) such that for every \(s \in S\) the higher cohomologies
  \(H^r(\PP^n_s, \F_s(m))\) vanish for all \(r \ge 1\), \(m \ge N\), and the base-change
  homomorphism \((\pi_* \F(m))|_s \to H^0(\PP^n_s, \F_s(m))\) is an isomorphism
  for \(m \ge N\), \(s \in S\).  Iterating the \(n=0\) stratification
  (\cref{lem:flat_locus_open}) on the coherent sheaves \(E_i := \pi_* \F(N+i)\)
  for \(i = 0, \dots, n\) produces a finite collection of locally-closed
  subschemes \(W_{e_0, \dots, e_n} \subseteq S\) indexed by tuples
  \((e_0, \dots, e_n) \in \Z^{n+1}\). Re-indexing each tuple by the unique
  numerical polynomial \(f \in \Q[\lambda]\) of degree \(\le n\) with
  \(f(N+i) = e_i\) gives a finite stratification \(\{W_f\}\) of \(S\), and the
  locally-closed subschemes \(S_f \subseteq W_f\) of
  \cref{thm:flattening_stratification_exists} are obtained by further
  intersecting with the closed-subscheme conditions that
  \(\pi_* \F(N+i)|_{W_f}\) be locally free of rank \(f(N+i)\) for all \(i \ge 0\)
  (terminating by Noetherianity).
\end{lemma}

\begin{proof}
  By \cref{lem:nonflat_locus_proper}, the family of Hilbert polynomials
  \(\{P_s\}\) on \(S\) is finite, and there is a finite stratification
  \(\{V_i\}\) of \(S\) on each of which \(\F\) is flat with locally-constant
  Hilbert polynomial.  Apply semi-continuity (cohomology and base change,
  \cite{hartshorne-ag}, III.12 / \cite{stacks-project}, tag 02KH for the \(H^0\)-version) on
  the noetherian \(V_i\) to obtain a uniform integer \(N_1\) with
  \(R^r \pi_* \F(m) = 0\) for \(r \ge 1\), \(m \ge N_1\), and \(H^r(\PP^n_s, \F_s(m)) = 0\)
  for all \(s \in S\), \(r \ge 1\), \(m \ge N_1\) (statement (B) of Nitsure's proof).
  Choose \(r_i \ge N_1\) on each \(V_i\) such that the base-change comparison
  \((\pi_* \F(m))|_{V_i} \to (\pi_i)_* \F_{V_i}(m)\) is an isomorphism for
  \(m \ge r_i\); composing with the cohomology-and-base-change isomorphism
  \(((\pi_i)_* \F_{V_i}(m))|_s \xrightarrow{\sim} H^0(\PP^n_s, \F_s(m))\) for
  the flat family on \(V_i\) yields uniform \(N = \max r_i\) such that
  \((\pi_* \F(m))|_s \to H^0(\PP^n_s, \F_s(m))\) is an isomorphism for \(m \ge N\)
  and all \(s \in S\) (statement (C)).

  Apply \cref{lem:flat_locus_open} sequentially to the coherent sheaves
  \(E_i = \pi_* \F(N+i)\) on \(S\) for \(i = 0, \dots, n\): \(E_0\) produces a
  stratification \(\{W_{e_0}\}\) of \(S\); on each \(W_{e_0}\) apply the lemma to
  \(E_1|_{W_{e_0}}\) to refine to \(\{W_{e_0, e_1}\}\), and continue. After \(n+1\)
  steps the \(W_{e_0, \dots, e_n}\) form a finite locally-closed stratification
  of \(S\) universal for the property ``each \(E_i\) pulls back to a locally-free
  \(\OO_T\)-module of constant rank \(e_i\)''.  Re-index \((e_0, \dots, e_n)\) by the
  unique numerical polynomial \(f\) of degree \(\le n\) with \(f(N+i) = e_i\) to get
  the family \(\{W_f\}\) indexed by polynomials.  At any \(s \in W_f\), statement
  (C) and vanishing (B) force \(P_s = f\) (since \(P_s\) has degree \(\le n\) and is
  determined by its values at \(N, N+1, \dots, N+n\)).

  The Hilbert-polynomial-coincidence locus \(|W_f|\) may have non-flat scheme
  structure; refine to \(S_f \subseteq W_f\) by intersecting with the
  closed-subscheme conditions ``\(\pi_* \F(N+i)|_{W_f}\) is locally-free of rank
  \(f(N+i)\)'' for all \(i \ge 0\).  These conditions are cut out by an increasing
  chain of coherent ideal sheaves \(I_0 \subseteq I_0 + I_1 \subseteq \cdots\) on
  \(W_f\); Noetherianity terminates the chain in finitely many steps, defining a
  coherent ideal sheaf \(I\) and a closed subscheme \(S_f \subseteq W_f\) with
  \(|S_f| = |W_f|\) on which \(\pi_* \F(N+i)\) is locally-free of rank \(f(N+i)\)
  for all \(i \ge 0\). The base-change-without-flatness lemma (cf.\ \cite{nitsure-hilbert-quot},
  \S3) yields some \(N' \ge N\) such that for \(i \ge N'\),
  \((\pi_* \F(i))|_{S_f} \to (\pi_{S_f})_* \F_{S_f}(i)\) is an isomorphism, hence
  the graded \(A\)-module \(\bigoplus_{i \ge N'} (\pi_{S_f})_* \F_{S_f}(i)\) is
  locally-free on \(S_f\), hence \(\F_{S_f}\) is flat over \(S_f\). The
  closure-of-strata property (iii) reduces to the corresponding property in
  the \(n=0\) case for \(\pi_* \F(p)\) on \(S\) at any \(p \ge N\) separating all
  Hilbert polynomials.
\end{proof}

\subsection{Consequences of the flattening construction}
\label{sec:flatstrat_lean_helpers}

The flattening construction also yields a countably indexed flat
stratification.  Further transfer lemmas give a direct Noetherian-induction
construction for proper morphisms and allow the resulting finite family to
be indexed injectively by integer-valued functions.

\begin{lemma}[Flat-locus stratification, special case \(n = 0\)]\leanok
  \label{lem:flat_locus_stratification_lean}
  \dcref{custom}
  \lean{AlgebraicGeometry.flatLocusStratification}
  \uses{def:coherent_sheaf_flat}
  Let \(S\) be a noetherian scheme and \(\F\) a coherent \(\OO_S\)-module.
  Then there is a family of immersions \(\iota_e : S_e \to S\) indexed by
  \(e \in \mathbb{N}\) whose images are pairwise disjoint and set-theoretically
  cover \(|S|\), and such that for each \(e\) the pullback of \(\F\) along
  \(\iota_e\) is flat over \(S_e\) in the sense of
  \cref{def:coherent_sheaf_flat}.  The index \(e\) is merely an enumeration;
  unlike the canonical strata of \cref{lem:flat_locus_open}, it is not
  required to equal \(\dim_{\kappa(s)} \F|_s\).
\end{lemma}

\begin{proof}
  \leanok
  \uses{thm:flattening_stratification_exists}
  Apply the existence theorem \cref{thm:flattening_stratification_exists} to
  the identity morphism \(\pi = \mathrm{id}_S\) and \(\F\): it yields a
  finite set \(I\), immersions \(\iota_f : V_f \to S\) with pairwise
  disjoint images covering \(|S|\), and flatness of the pullback of \(\F\)
  to \(S \times_S V_f\) over \(V_f\).  The second projection
  \(p_f : S \times_S V_f \to V_f\) of the fibre product along the identity
  is an isomorphism, and the first projection factors as
  \(\mathrm{pr}_1 = p_f \circ \iota_f\); take \(S \times_S V_f\) itself
  as the stratum with \(\mathrm{pr}_1\) as its immersion, so disjointness
  and covering of images transfer along the surjection \(p_f\).  Flatness of
  the pullback over the stratum itself follows because the predicate of
  \cref{def:coherent_sheaf_flat} is insensitive to composing the base leg
  with an isomorphism: an affine open of the target corresponds under the
  isomorphism to an affine open of the source, the two section-ring actions
  on an unchanged section module differ by composition with a bijective ring
  map, and flatness passes along such a factorization since the target ring
  is a free module of rank one over the source.  Finally choose a bijection
  of \(I\) with an initial segment of \(\mathbb{N}\) and pad the family at
  the remaining indices by the empty scheme: its section rings are trivial,
  so every section module over them is flat, and the empty image preserves
  disjointness and covering.
\end{proof}

\subsection{Flatness transfer to strata}
\label{subsec:flat_transfer_strata}

The Noetherian-induction reduction is assembled from four transfer lemmas:
flatness under pushout base change, the affine-piece analysis of a fibre
product, affine-locality of flatness over the base, and the existence of a
flat stratum over one irreducible component.

\begin{lemma}[Flatness under pushout base change]

  \leanok
  \label{lem:flat_pushout_base_change}
  \dcref{custom}
  \lean{Module.Flat.of_isPushout}
  Let \(R \to S\) and \(R \to A\) be commutative ring maps with pushout
  \(B = S \otimes_R A\), let \(M\) be an \(A\)-module which is flat over
  \(R\), and let \(N\) be a \(B\)-module which is a base change of \(M\)
  along \(A \to B\), i.e.\ \(N \cong B \otimes_A M\) as \(B\)-modules.  Then
  \(N\) is flat over \(S\).
\end{lemma}

\begin{proof}
  \leanok
  It suffices to prove that \(B \otimes_A M\) is flat over \(S\), and for
  this we show that the canonical \(R\)-linear map
  \(\eta : M \to B \otimes_A M\), \(m \mapsto 1 \otimes m\), exhibits
  \(B \otimes_A M\) as the base change \(S \otimes_R M\); flatness over
  \(S\) then follows from flatness of \(M\) over \(R\) since base change
  along \(R \to S\) preserves flatness.  We verify the universal property:
  let \(Q\) be an \(S\)-module and \(g : M \to Q\) an \(R\)-linear map.  The
  family \(\beta : A \to \operatorname{Hom}_R(M, Q)\),
  \(\beta(a)(m) = g(am)\), is \(R\)-linear into an \(S\)-module, so by the
  universal property of the ring base change \(B = S \otimes_R A\) it lifts
  to an \(S\)-linear map \(\Lambda : B \to \operatorname{Hom}_R(M, Q)\) with
  \(\Lambda(1 \otimes a) = \beta(a)\).  The pairing
  \((b, m) \mapsto \Lambda(b)(m)\) is additive, \(S\)-linear in \(b\), and
  \(A\)-balanced: on generators \(b = s (1 \otimes a')\) of \(B\) one has
  \(\Lambda((1 \otimes a) b)(m) = s\, g(a a' m) = \Lambda(b)(a m)\), and the
  general case follows since such elements generate \(B\) additively.  Hence
  the pairing descends to an \(S\)-linear map
  \(h : B \otimes_A M \to Q\) with \(h \circ \eta = g\).  Uniqueness: any
  \(S\)-linear \(h'\) with \(h' \circ \eta = g\) agrees with \(h\) on
  elements \((s(1 \otimes a)) \otimes m = s\,(1 \otimes am)\), hence on all
  of \(B \otimes_A M\).
\end{proof}

\begin{lemma}[Flat sections on an affine piece of a fibre product]

\leanok
  \label{lem:flat_section_pullback_piece}
  \dcref{custom}
  \lean{AlgebraicGeometry.flat_section_pullback_piece}
  \uses{def:coherent_sheaf_flat, lem:flat_pushout_base_change,
    def:pullback_app_isoTensor_sigma}
  Let
  \[
    \begin{array}{ccc}
      Y & \xrightarrow{\;g\;} & X \\
      \downarrow{\scriptstyle i_Y} & & \downarrow{\scriptstyle i_X} \\
      T & \xrightarrow{\;f\;} & S
    \end{array}
  \]
  be a cartesian square of schemes, and let \(U_S \subseteq S\),
  \(U_X \subseteq i_X^{-1}(U_S)\), \(U_T \subseteq f^{-1}(U_S)\) be affine
  opens.  Consider the \emph{piece}
  \(U_Y := g^{-1}(U_X) \cap i_Y^{-1}(U_T) \subseteq Y\).  Then \(U_Y\) is an
  affine open, the section-ring square
  \(\Gamma(S, U_S) \to \Gamma(X, U_X),\ \Gamma(S, U_S) \to \Gamma(T, U_T)\)
  with corner \(\Gamma(Y, U_Y)\) is a pushout of commutative rings, and for
  every quasi-coherent \(\OO_X\)-module \(\F\) such that
  \(\Gamma(\F, U_X)\) is flat over \(\Gamma(S, U_S)\), the sections
  \(\Gamma(g^* \F, U_Y)\) are flat over \(\Gamma(T, U_T)\).  If moreover
  \(\Gamma(\F, U_X)\) is a finite \(\Gamma(X, U_X)\)-module, then
  \(\Gamma(g^* \F, U_Y)\) is a finite \(\Gamma(Y, U_Y)\)-module.
\end{lemma}

\begin{proof}
  \leanok
  Restricting the cartesian square over the affine cospan
  \(U_X \to U_S \leftarrow U_T\) exhibits \(U_Y\) as the fibre product
  \(U_X \times_{U_S} U_T\) of affine schemes, which is affine.  Applying
  global sections to this pullback and using the anti-equivalence between
  affine schemes and commutative rings yields the pushout of section rings.
  Since \(U_Y \subseteq g^{-1}(U_X)\) with \(U_X\), \(U_Y\) affine and
  \(\F\) quasi-coherent, the base-change section formula
  (\cref{def:pullback_app_isoTensor_sigma}, Stacks 01HQ/01I8) identifies
  \(\Gamma(g^* \F, U_Y) \cong
    \Gamma(Y, U_Y) \otimes_{\Gamma(X, U_X)} \Gamma(\F, U_X)\)
  as \(\Gamma(Y, U_Y)\)-modules, i.e.\ exhibits \(\Gamma(g^*\F, U_Y)\) as a
  base change of \(\Gamma(\F, U_X)\) along
  \(\Gamma(X, U_X) \to \Gamma(Y, U_Y)\).  Flatness over \(\Gamma(T, U_T)\)
  is now \cref{lem:flat_pushout_base_change} applied to the section-ring
  pushout square; finiteness is preserved because a base change of a finite
  module is finite.
\end{proof}

\begin{lemma}[Flatness over the base is affine-local]

\leanok
  \label{lem:flat_affine_local}
  \dcref{custom}
  \lean{AlgebraicGeometry.flat_section_of_affine_cover}
  \uses{def:coherent_sheaf_flat, lem:appLE_base_res_apply}
  Let \(q : Y \to T\) be a morphism of schemes and \(\mathcal G\) a quasi-coherent
  \(\OO_Y\)-module.  Suppose given affine opens
  \(U_j \subseteq T\) and \(W_j \subseteq q^{-1}(U_j)\) (\(j \in J\)) such
  that the \(W_j\) cover \(Y\) and each \(\Gamma(\mathcal G, W_j)\) is flat over
  \(\Gamma(T, U_j)\).  Then for \emph{every} pair of affine opens
  \(U \subseteq T\), \(W \subseteq q^{-1}(U)\), the sections
  \(\Gamma(\mathcal G, W)\) are flat over \(\Gamma(T, U)\).
\end{lemma}

\begin{proof}
  \leanok
  This is the two-layer basic-open reduction of
  \cref{lem:flat_section_pair}, with the freeness input replaced by chart
  flatness.  Fix an affine pair \((U, W)\).  For each \(x \in W\) choose a
  chart \(W_j \ni x\); since \(q(x) \in U \cap U_j\) and both are affine,
  there is a basic open \(D(a) = D(a')\) of \(T\) common to \(U\) and
  \(U_j\) containing \(q(x)\), and then a basic open
  \(D(b) = D(c) \ni x\) common to \(W\) and \(W_j\) with
  \(D(b) \subseteq q^{-1}(D(a))\).  The elements \(b\) generate the unit
  ideal of \(\Gamma(Y, W)\), and each restriction
  \(\Gamma(\mathcal G, W) \to \Gamma(\mathcal G, D(b))\) is a localization at the powers of
  \(b\) (\cref{lem:qcoh_section_localization_basicOpen}), so flatness of
  \(\Gamma(\mathcal G, W)\) over \(\Gamma(T, U)\) may be checked on the pieces
  \(\Gamma(\mathcal G, D(b))\).  On a piece: flatness of \(\Gamma(\mathcal G, W_j)\) over
  \(\Gamma(T, U_j)\) passes to the fibre-side localization
  \(\Gamma(\mathcal G, D(c))\) (mixed-base stability of flatness under localization
  of the module at a submonoid of the fibre algebra), and then the base ring
  is exchanged twice along
  \(\Gamma(T, U_j) \to \Gamma(T, D(a')) = \Gamma(T, D(a))\) and
  \(\Gamma(T, U) \to \Gamma(T, D(a))\) (for a module over a localized ring,
  flatness over the ring and over its localization agree; the compatibility
  of the induced module structures is \cref{lem:appLE_base_res_apply}).
\end{proof}

\begin{lemma}[Reduced subscheme of an irreducible closed subset is integral]

\leanok
  \label{lem:vanishing_ideal_subscheme_integral}
  \dcref{custom}
  \lean{AlgebraicGeometry.isIntegral_vanishingIdeal_subscheme}
  Let \(S\) be a scheme and \(Z \subseteq S\) a closed subset.  The
  subscheme cut out by the vanishing ideal sheaf of \(Z\) (the reduced
  induced structure) is reduced; if \(Z\) is irreducible (in particular
  non-empty), it is integral.
\end{lemma}

\begin{proof}
  \leanok
  The vanishing ideal sheaf is radical: it equals the vanishing ideal of
  its own support (Galois connection between supports and vanishing
  ideals), which is the radical.  The subscheme is covered by the affine
  charts \(\Spec(\Gamma(S, U)/I(U))\) for \(U\) ranging over affine opens
  of \(S\), and each \(I(U)\) is a radical ideal, so each chart is the
  spectrum of a reduced ring; hence the subscheme is reduced.  Its
  underlying space is the subspace \(Z\), which is irreducible by
  hypothesis, and a reduced scheme with irreducible non-empty underlying
  space is integral.
\end{proof}

\begin{lemma}[Flat stratum over an irreducible component]

\leanok
  \label{lem:flat_stratum_component}
  \dcref{custom}
  \lean{AlgebraicGeometry.flat_stratum_of_irreducible}
  \uses{def:coherent_sheaf_flat, thm:generic_flatness,
    lem:vanishing_ideal_subscheme_integral, lem:flat_section_pullback_piece,
    lem:flat_affine_local}
  Let \(S\) be a noetherian scheme, \(\pi : X \to S\) proper, \(\F\) a
  coherent \(\OO_X\)-module, and \(Z_c \subseteq S\) an irreducible closed
  subset with reduced subscheme \(T \to S\).  Then there is a non-empty open
  \(\Omega \subseteq T\) such that for \emph{every} open
  \(\mathit{St} \subseteq \Omega\), regarded as a locally closed subscheme
  \(\mathit{St} \to S\), the pullback of \(\F\) to
  \(X \times_S \mathit{St}\) is flat over \(\mathit{St}\) in the sense of
  \cref{def:coherent_sheaf_flat}.
\end{lemma}

\begin{proof}
  \leanok
  \(T\) is integral (\cref{lem:vanishing_ideal_subscheme_integral}) and
  noetherian (a closed subscheme of a noetherian scheme), and the base
  change \(p : X \times_S T \to T\) is proper, hence quasi-compact and
  locally of finite type.  The pullback of \(\F\) to \(X \times_S T\) is
  quasi-coherent, and it admits affine charts with finite section modules:
  around any point, choose affine opens \(U_S \subseteq S\),
  \(U_T \subseteq T\) above it, and an affine chart \(W_X \subseteq X\)
  with finite sections (Stacks 01PC); the corresponding piece is an affine
  chart of \(X \times_S T\) whose sections are the base change of
  \(\Gamma(\F, W_X)\), hence finite (\cref{lem:flat_section_pullback_piece}),
  and this finiteness localizes to a basic open inside any given
  neighbourhood.  Generic flatness (\cref{thm:generic_flatness}, in its
  quasi-coherent form with this chart supply) yields a non-empty open
  \(\Omega \subseteq T\) with flat sections on all affine pairs below
  \(\Omega\).

  Now let \(\mathit{St} \subseteq \Omega\) be open.  The stratum family
  \(X \times_S \mathit{St} \to \mathit{St}\) maps to the component family by
  the comparison morphism
  \(\kappa : X \times_S \mathit{St} \to X \times_S T\), and the square
  formed by \(\kappa\) and the two projections over
  \(\mathit{St} \hookrightarrow T\) is cartesian (cancellation of pasted
  pullback squares).  Cover \(X \times_S \mathit{St}\) by pieces
  \(g^{-1}(W_X) \cap q^{-1}(U')\) with \(U' \subseteq \mathit{St}\) affine
  lying over an affine \(U_1 \subseteq \Omega\); on each such piece,
  applying \cref{lem:flat_section_pullback_piece} to the cartesian
  \(\kappa\)-square with flatness input from the generic-flatness output on
  the corresponding piece of \(X \times_S T\) (an affine open below
  \(\Omega\)) yields flat sections over \(U'\).  By affine-locality
  (\cref{lem:flat_affine_local}) all affine pairs of the stratum family
  have flat sections; finally the pulled-back sheaf along \(\kappa\)
  composed with the projection is identified with the pullback of \(\F\)
  along the projection of the stratum family (functoriality of pullback,
  \(g \circ \kappa = g'\)), so the flatness predicate transports along this
  isomorphism of module sheaves.
\end{proof}

\begin{lemma}[Noetherian-induction reduction to the flat case]\leanok
  \label{lem:flat_locus_reduction_lean}
  \dcref{custom}
  \lean{AlgebraicGeometry.flatLocusReduction}
  \uses{def:coherent_sheaf_flat, thm:generic_flatness}
  Let \(S\) be a noetherian scheme, \(\pi : X \to S\) a proper morphism,
  and \(\F\) a coherent \(\OO_X\)-module. Then there exist finitely many
  immersions \(\iota_i : V_i \to S\) with pairwise disjoint images
  set-theoretically covering \(|S|\), such that for each \(i\) the pullback of
  \(\F\) to \(X \times_S V_i\) is flat over \(V_i\) in the sense of
  \cref{def:coherent_sheaf_flat}.  The strata are obtained by successively
  removing non-empty flat open patches supplied by generic flatness
  (\cref{thm:generic_flatness}).

  Noetherianity of \(S\) (rather than mere local noetherianity) is necessary
  for the \emph{finiteness} of the family: on
  \(S = \coprod_{n \ge 1} \mathbb{A}^n\), equip the \(n\)-th component with
  the coherent sheaf \(\bigoplus_{k \le n} (i_{V_k})_* \OO_{V_k}\) for a
  strictly nested flag \(V_1 \supset \cdots \supset V_n\) of integral closed
  subvarieties.  A locally-closed stratum on which this sheaf is flat (hence
  locally free) meets the flag levels in relatively clopen subsets, so a
  stratum through the generic point of \(V_k\) cannot also contain the
  generic point of \(V_{k'}\) for \(k' \neq k\): each neighbourhood of the
  latter meets the dense lower level.  The \(n\)-th component therefore
  requires at least \(n + 1\) strata, and no finite family covers all
  components.
\end{lemma}

\begin{proof}
  \leanok
  \uses{lem:nonflat_locus_proper, thm:generic_flatness,
    lem:flat_stratum_component, lem:vanishing_ideal_subscheme_integral}
  By well-founded induction on the closed subsets of \(S\) (the noetherian
  hypothesis makes strict inclusion of closed subsets well-founded), we prove:
  for every closed \(Z \subseteq S\) there is a finite family of immersions
  with pairwise disjoint images whose union is exactly \(Z\) and with flat
  pullback of \(\F\) on each member.  For \(Z = \emptyset\) take the empty
  family.

  For \(Z \neq \emptyset\), let \(Z_1, \dots, Z_r \subseteq S\) be the images
  of the (finitely many, by noetherianity) irreducible components of the
  subspace \(Z\); each \(Z_c\) is an irreducible closed subset of \(S\)
  contained in \(Z\), and together they cover \(Z\).
  \Cref{lem:flat_stratum_component} applied to the reduced subscheme
  \(T_c \to S\) of \(Z_c\) yields a non-empty open \(\Omega_c \subseteq T_c\)
  all of whose open subsets are flat strata.  Let
  \(E_c \subseteq S\) be the open complement of the union of the other
  components, and take as stratum the open
  \(\mathit{St}_c := \Omega_c \cap (T_c \to S)^{-1}(E_c) \subseteq T_c\),
  a locally closed subscheme of \(S\) via
  \(\mathit{St}_c \hookrightarrow T_c \to S\), with flat pullback by the
  choice of \(\Omega_c\).

  The images of the \(\mathit{St}_c\) are pairwise disjoint: the image of
  \(\mathit{St}_c\) is contained in \(Z_c \setminus \bigcup_{c' \neq c} Z_{c'}\).
  Their union \(V\) is relatively open in \(Z\): realizing the open image of
  each \(\mathit{St}_c\) inside \(Z_c\) by an open \(W_c \subseteq S\), one has
  \(V = Z \cap \bigcup_c (W_c \cap E_c)\), because a point of
  \(Z \cap W_c \cap E_c\) avoids all components other than \(Z_c\), hence lies
  in \(Z_c\), hence in the image of \(\mathit{St}_c\).  Moreover \(V\) is
  non-empty: for any component \(Z_{c}\), irredundancy of the component
  decomposition shows \(Z_c \not\subseteq \bigcup_{c' \neq c} Z_{c'}\), so the
  open \((T_c \to S)^{-1}(E_c)\) is non-empty, and two non-empty opens of the
  irreducible space \(T_c\) intersect, so \(\mathit{St}_c \neq \emptyset\).

  Therefore \(Z' := Z \setminus V\) is a closed subset strictly smaller than
  \(Z\); the induction hypothesis provides a finite disjoint flat family with
  union \(Z'\), and adjoining the strata \(\mathit{St}_c\) gives the required
  family for \(Z\) (disjointness across the two groups holds because the
  recursive family lies in \(Z'\), which is disjoint from \(V\)).  Applying
  the statement to \(Z = S\) proves the lemma.
\end{proof}

\begin{lemma}[Assembly with injectively labelled strata]\leanok
  \label{lem:flat_locus_assembly_lean}
  \dcref{custom}
  \lean{AlgebraicGeometry.flatLocusAssembly}
  \uses{def:coherent_sheaf_flat}
  Let \(S\) be a noetherian scheme, \(\pi : X \to S\) a proper morphism,
  and \(\F\) a coherent \(\OO_X\)-module. Then there exist a finite index set
  \(I\), immersions \(\iota_f : S_f \to S\) (\(f \in I\)) with pairwise disjoint
  images set-theoretically covering \(|S|\), and an injective label map
  \(P : I \to (\mathbb{N} \to \Z)\) (so \(f = g \iff P_f = P_g\)), such that for
  each \(f\) the pullback of \(\F\) to \(X \times_S S_f\) is flat over \(S_f\) in
  the sense of \cref{def:coherent_sheaf_flat}.
\end{lemma}

\begin{proof}
  \leanok
  \uses{lem:flat_locus_reduction_lean}
  Since the statement requires of \(P\) only injectivity, it follows from the
  reduction \cref{lem:flat_locus_reduction_lean}: take its stratification and
  label the finitely many strata by any injection
  \(I \hookrightarrow \mathbb{N} \to \Z\) (for instance constant functions
  with pairwise distinct values).  Requiring these labels to be the fibrewise
  Hilbert polynomials gives the refinement of
  \cref{lem:noetherian_induction_strata}.
\end{proof}

\section{The entry ideal of a matrix presentation}
\label{sec:entry_ideal}

The rank strata of \cref{lem:flat_locus_open} carry a canonical --- in
general non-reduced --- scheme structure, cut out locally by the ideal of
entries of the relation matrix of a presentation.  This section develops the
commutative-algebra substrate: presentations of a module by \(e\)
generators, the entry ideal, its base-change behaviour, Nitsure's local
universal property, independence of the chosen presentation, and the
Nakayama prolongation producing \(e\)-generator presentations near a point
of fiber rank \(e\).  Throughout, \(R\) is a commutative ring and \(M\) an
\(R\)-module.

\begin{definition}[Matrix presentation]
  \label{def:matrix_presentation}
  \lean{Module.MatrixPresentation}
  \leanok
  \dcref{custom}
  A \emph{matrix presentation} of \(M\) by \(e\) generators and \(m\)
  relations consists of a matrix
  \(\psi \in \operatorname{Mat}_{e \times m}(R)\) and a surjection
  \(\pi : R^{\oplus e} \to M\) making the sequence
  \[
    R^{\oplus m} \xrightarrow{\ \psi\ } R^{\oplus e}
      \xrightarrow{\ \pi\ } M \to 0
  \]
  exact.
\end{definition}

\begin{lemma}[Finite presentation]
  \label{lem:matrix_presentation_fp}
  \dcref{custom}
  \lean{Module.MatrixPresentation.finitePresentation}
  \leanok
  \uses{def:matrix_presentation}
  A module admitting a matrix presentation is finitely presented.
\end{lemma}

\begin{proof}
  \leanok
  \(\pi\) is a surjection from a finite free module whose kernel, the image
  of \(\psi\), is generated by the \(m\) columns of \(\psi\), hence finitely
  generated.
\end{proof}

\begin{definition}[Entry ideal]
  \label{def:entry_ideal}
  \lean{Module.MatrixPresentation.entryIdeal}
  \leanok
  \uses{def:matrix_presentation}
  \dcref{custom}
  The \emph{entry ideal} \(I(\psi) \subseteq R\) of a matrix presentation is
  the ideal generated by the \(e \cdot m\) entries \(\psi_{ij}\) of the
  relation matrix.  It is the ideal of \(1 \times 1\) minors of \(\psi\);
  for an \(e\)-generator presentation this is the Fitting ideal
  \(\operatorname{Fitt}_{e-1}(M)\).
\end{definition}

\begin{lemma}[Vanishing criterion]
  \label{lem:entry_ideal_vanish}
  \dcref{custom}
  \lean{Module.MatrixPresentation.relMatrix_map_eq_zero_iff}
  \leanok
  \uses{def:entry_ideal}
  For an \(R\)-algebra \(A\), the entrywise image matrix
  \(\psi \otimes A \in \operatorname{Mat}_{e \times m}(A)\) vanishes if and
  only if \(I(\psi) \subseteq \ker(R \to A)\).
\end{lemma}

\begin{proof}
  \leanok
  Both sides assert that every entry of \(\psi\) maps to \(0\) in \(A\).
\end{proof}

\begin{definition}[Base change of a presentation]
  \label{def:matrix_presentation_base_change}
  \lean{Module.MatrixPresentation.baseChange}
  \leanok
  \uses{def:matrix_presentation}
  \dcref{custom}
  Let \(A\) be an \(R\)-algebra.  The matrix \(\psi \otimes A\) obtained by
  applying \(R \to A\) entrywise, together with the tensored projection
  \(A^{\oplus e} \cong A \otimes_R R^{\oplus e}
    \xrightarrow{1 \otimes \pi} A \otimes_R M\),
  is a matrix presentation of \(A \otimes_R M\) by \(e\) generators and
  \(m\) relations: the tensored sequence
  \(A^{\oplus m} \to A^{\oplus e} \to A \otimes_R M \to 0\) remains exact by
  right-exactness of the tensor product.
\end{definition}

\begin{lemma}[Entry ideal under base change]
  \label{lem:entry_ideal_base_change}
  \dcref{custom}
  \lean{Module.MatrixPresentation.entryIdeal_baseChange}
  \leanok
  \uses{def:entry_ideal, def:matrix_presentation_base_change}
  The entry ideal of the base-changed presentation is the image ideal:
  \(I(\psi \otimes A) = I(\psi) \cdot A\).
\end{lemma}

\begin{proof}
  \leanok
  The generators of \(I(\psi \otimes A)\) are exactly the images of the
  generators of \(I(\psi)\).
\end{proof}

\begin{lemma}[Vanishing relations give an isomorphism]
  \label{lem:presentation_iso_of_vanish}
  \dcref{custom}
  \lean{Module.MatrixPresentation.projEquiv}
  \leanok
  \uses{def:matrix_presentation}
  If \(\psi = 0\), then \(\pi : R^{\oplus e} \to M\) is an isomorphism.
\end{lemma}

\begin{proof}
  \leanok
  \(\ker \pi = \operatorname{im} \psi = 0\) by exactness, and \(\pi\) is
  surjective.
\end{proof}

\begin{lemma}[Vanishing relations give a free module]
  \label{lem:presentation_free_of_vanish}
  \dcref{custom}
  \lean{Module.MatrixPresentation.free_of_relMatrix_eq_zero}
  \leanok
  \uses{lem:presentation_iso_of_vanish}
  If \(\psi = 0\), then \(M\) is free.
\end{lemma}

\begin{proof}
  \leanok
  Transport freeness along the isomorphism of
  \cref{lem:presentation_iso_of_vanish}.
\end{proof}

\begin{lemma}[Vanishing relations give constant rank]
  \label{lem:presentation_rank_of_vanish}
  \dcref{custom}
  \lean{Module.MatrixPresentation.rankAtStalk_eq_of_relMatrix_eq_zero}
  \leanok
  \uses{lem:presentation_iso_of_vanish}
  If \(\psi = 0\), then the localization \(M_{\mathfrak p}\) is free of rank
  \(e\) over \(R_{\mathfrak p}\) for every prime
  \(\mathfrak p \subseteq R\).
\end{lemma}

\begin{proof}
  \leanok
  By \cref{lem:presentation_iso_of_vanish}, \(M \cong R^{\oplus e}\); a ring
  with a prime ideal is nontrivial, so \(R^{\oplus e}\) has constant rank
  \(e\) at every prime.
\end{proof}

\begin{theorem}[Flat everywhere-rank-\(e\) modules have vanishing relation
  matrix]
  \label{thm:entry_ideal_flat_vanish}
  \lean{Module.MatrixPresentation.relMatrix_eq_zero_of_flat}
  \leanok
  \uses{def:matrix_presentation}
  \dcref{custom}
  Let \(M\) carry a matrix presentation \((\psi, \pi)\) by \(e\) generators.
  If \(M\) is flat and \(M_{\mathfrak p}\) has rank \(e\) over
  \(R_{\mathfrak p}\) for every prime \(\mathfrak p\), then \(\psi = 0\).
  Equivalently, a surjection \(R^{\oplus e} \twoheadrightarrow M\) onto a
  locally free module of everywhere-rank \(e\) has zero kernel.
\end{theorem}

\begin{proof}
  \leanok
  \uses{lem:matrix_presentation_fp}
  By \cref{lem:matrix_presentation_fp}, \(M\) is finitely presented, hence
  --- being flat --- projective.  The surjection \(\pi\) therefore splits:
  choose \(s : M \to R^{\oplus e}\) with \(\pi \circ s = \mathrm{id}\).
  The splitting yields an isomorphism
  \(R^{\oplus e} \cong M \times \ker \pi\), \(x \mapsto
  (\pi x,\, x - s(\pi x))\), so \(K := \ker \pi\) is a direct summand of a
  finite free module, hence itself finite projective (in particular flat).
  Localizing at any prime \(\mathfrak p\) and counting ranks,
  \(e = \operatorname{rk}_{\mathfrak p} R^{\oplus e}
     = \operatorname{rk}_{\mathfrak p} M + \operatorname{rk}_{\mathfrak p} K
     = e + \operatorname{rk}_{\mathfrak p} K\),
  so \(K\) has rank \(0\) at every prime; a finite flat module of
  everywhere-rank \(0\) is zero.  Hence
  \(\operatorname{im} \psi = \ker \pi = 0\) by exactness, and every entry
  \(\psi_{ij} = (\psi(\delta_j))_i\) vanishes.
\end{proof}

\begin{lemma}[Comparison of entry ideals]
  \label{lem:entry_ideal_le}
  \dcref{custom}
  \lean{Module.MatrixPresentation.entryIdeal_le_entryIdeal}
  \leanok
  \uses{def:entry_ideal}
  If \((\psi, \pi)\) and \((\psi', \pi')\) are two matrix presentations of
  \(M\) by \(e\) generators (with possibly different numbers of relations),
  then \(I(\psi') \subseteq I(\psi)\).
\end{lemma}

\begin{proof}
  \leanok
  \uses{def:matrix_presentation_base_change, lem:entry_ideal_vanish,
    lem:presentation_free_of_vanish, lem:presentation_rank_of_vanish,
    thm:entry_ideal_flat_vanish}
  Base change both presentations to \(A := R / I(\psi)\).  The entries of
  \(\psi \otimes A\) vanish by construction
  (\cref{lem:entry_ideal_vanish}), so by
  \cref{lem:presentation_free_of_vanish} and \cref{lem:presentation_rank_of_vanish}
  the module \(A \otimes_R M\) is free (hence flat) of rank \(e\) at every
  prime of \(A\).  Applying \cref{thm:entry_ideal_flat_vanish} to the
  base-changed presentation \((\psi' \otimes A, \pi' \otimes A)\) gives
  \(\psi' \otimes A = 0\), i.e.\ every entry of \(\psi'\) lies in
  \(\ker(R \to R/I(\psi)) = I(\psi)\).
\end{proof}

\begin{theorem}[Presentation-independence of the entry ideal]
  \label{thm:entry_ideal_presentation_independent}
  \lean{Module.MatrixPresentation.entryIdeal_eq}
  \leanok
  \uses{def:entry_ideal}
  \dcref{custom}
  Any two matrix presentations of \(M\) by \(e\) generators have the same
  entry ideal.  (This is the presentation-independence of the Fitting ideal
  \(\operatorname{Fitt}_{e-1}\), and is what lets the local entry-ideal
  subschemes of \cref{lem:flat_locus_open} glue.)
\end{theorem}

\begin{proof}
  \leanok
  \uses{lem:entry_ideal_le}
  Apply \cref{lem:entry_ideal_le} in both directions.
\end{proof}

\begin{definition}[Fiber rank]
  \label{def:fiber_rank}
  \lean{Ideal.fiberRank}
  \leanok
  \dcref{custom}
  For a prime \(\mathfrak p \subseteq R\) with residue field
  \(\kappa(\mathfrak p)\), the \emph{fiber rank} of \(M\) at
  \(\mathfrak p\) is
  \(\dim_{\kappa(\mathfrak p)} \bigl(\kappa(\mathfrak p) \otimes_R M\bigr)\).
\end{definition}

\begin{theorem}[Point-set of the entry-ideal locus]
  \label{thm:entry_ideal_prime_iff}
  \lean{Module.MatrixPresentation.entryIdeal_le_prime_iff}
  \leanok
  \uses{def:entry_ideal, def:fiber_rank}
  \dcref{custom}
  Let \((\psi, \pi)\) be a matrix presentation of \(M\) by \(e\) generators
  and \(\mathfrak p \subseteq R\) a prime.  Then
  \(I(\psi) \subseteq \mathfrak p\) if and only if the fiber rank of \(M\)
  at \(\mathfrak p\) equals \(e\).
\end{theorem}

\begin{proof}
  \leanok
  \uses{def:matrix_presentation_base_change, lem:entry_ideal_vanish,
    lem:presentation_iso_of_vanish, thm:entry_ideal_flat_vanish}
  Base change the presentation to the residue field
  \(\kappa(\mathfrak p)\), whose structure map has kernel exactly
  \(\mathfrak p\).  If \(I(\psi) \subseteq \mathfrak p\), the fiber
  presentation has vanishing relation matrix
  (\cref{lem:entry_ideal_vanish}), so the fiber is isomorphic to
  \(\kappa(\mathfrak p)^{\oplus e}\)
  (\cref{lem:presentation_iso_of_vanish}) and has dimension \(e\).
  Conversely, if the fiber has dimension \(e\), then --- every module over
  a field being free, and free modules over a field having rank equal to
  their dimension at the unique relevant localization --- the fiber
  presentation satisfies the hypotheses of
  \cref{thm:entry_ideal_flat_vanish}, so its relation matrix vanishes and
  every entry of \(\psi\) lies in
  \(\ker(R \to \kappa(\mathfrak p)) = \mathfrak p\).
\end{proof}

\begin{theorem}[Nakayama prolongation of a fiber basis]
  \label{thm:nakayama_prolongation}
  \lean{Module.FinitePresentation.exists_matrixPresentation_of_isLocalizedModule}
  \leanok
  \uses{def:matrix_presentation, def:fiber_rank}
  \dcref{custom}
  Let \(M\) be finitely presented over \(R\), let \(\mathfrak p \subseteq R\)
  be a prime, and let \(e\) be the fiber rank of \(M\) at \(\mathfrak p\).
  Then there exists \(g \notin \mathfrak p\) such that every localization of
  \(M\) away from \(g\) --- that is, every module \(N\) over an \(R\)-algebra
  \(R'\) realizing the localizations of \(R\) and \(M\) at the powers of
  \(g\) --- admits a matrix presentation by exactly \(e\) generators.
\end{theorem}

\begin{proof}
  \leanok
  \uses{lem:matrix_presentation_fp}
  Every element of the fiber \(\kappa(\mathfrak p) \otimes_R M\) becomes,
  after scaling by some \(r \notin \mathfrak p\) (a unit of
  \(\kappa(\mathfrak p)\)), a pure tensor \(1 \otimes s\) with \(s \in M\):
  this holds for pure tensors \(c \otimes s\) because
  \(\kappa(\mathfrak p)\) is a localization of \(R/\mathfrak p\)-fractions
  (clear the denominator of \(c\)), and sums are handled by scaling both
  summands by the product of their two scalars.  Scaling each vector of a
  basis of the fiber this way produces \(m_1, \dots, m_e \in M\) whose
  images \(1 \otimes m_i\) again form a basis of the fiber (they differ from
  the original basis by unit scalars).

  Let \(\alpha : R^{\oplus e} \to M\) send the standard basis to the
  \(m_i\), and let \(Q := M / \operatorname{im} \alpha\); it is finitely
  generated.  Since the \(1 \otimes m_i\) span the fiber and the tensor
  product is right exact, \(\kappa(\mathfrak p) \otimes_R Q = 0\); as
  \(\kappa(\mathfrak p) \otimes_R Q\) is the fiber of the finite
  \(R_{\mathfrak p}\)-module \(Q_{\mathfrak p}\) at the maximal ideal,
  Nakayama gives \(Q_{\mathfrak p} = 0\).  Hence \(\mathfrak p\) lies
  outside the support of the finite module \(Q\), which is the zero locus
  of its annihilator, so there is \(g \in \operatorname{Ann}(Q)\) with
  \(g \notin \mathfrak p\); thus \(g \cdot M \subseteq
  \operatorname{im} \alpha\).

  Now let \(R'\) and \(N\) realize the localizations of \(R\) and \(M\)
  away from \(g\), with structure map \(l : M \to N\).  The images
  \(l(m_i)\) generate \(N\) over \(R'\): any \(n \in N\) satisfies
  \(g^k n = l(x)\) for some \(x \in M\) and some \(k\), the element
  \(g x\) lies in \(\operatorname{im} \alpha\), and \(g\) and \(g^k\) are
  units in \(R'\).  This gives a surjection
  \(\alpha' : R'^{\oplus e} \to N\).  Localization is a base change, so
  \(N\) is finitely presented over \(R'\)
  (\cref{lem:matrix_presentation_fp} for \(M\), transported along the base
  change); hence \(\ker \alpha'\) is finitely generated, and any finite
  generating family of \(\ker \alpha'\), arranged as the columns of a
  matrix, completes \(\alpha'\) to a matrix presentation of \(N\) by \(e\)
  generators.
\end{proof}

\section{Rank strata as closed subschemes}
\label{sec:flatstrat_rank_stratum}

The entry-ideal commutative algebra of \cref{sec:entry_ideal} is local to an
affine chart.  This section globalises it: for a quasi-coherent module
\(\mathcal G\) on a scheme \(X\) it assembles the local entry ideals into a
quasi-coherent ideal sheaf whose closed subscheme --- the \emph{rank-\(e\)
stratum} --- has support exactly the locus where the fibre of \(\mathcal G\)
has dimension \(e\).  This is the geometric substrate of Nitsure's
\(n = 0\) construction (\cref{lem:flat_locus_open}).  Throughout \(X\) is a
scheme and \(\mathcal G\) a quasi-coherent \(\OO_X\)-module; \(\Gamma(\mathcal G, U)\)
denotes its sections over an open \(U\), a module over \(\Gamma(X, U) = \OO_X(U)\).

\subsection{Base change of the fibre rank}

\begin{theorem}[Base-change invariance of the fibre rank]
  \label{thm:fiber_rank_base_change}
  \dcref{custom}
  \lean{Ideal.fiberRank_baseChange}
  \leanok
  \uses{def:fiber_rank}
  Let \(A\) be an \(R\)-algebra, \(M\) an \(R\)-module, and \(\mathfrak q\) a
  prime of \(A\) with contraction \(\mathfrak p = \mathfrak q \cap R\).  Then
  the fibre rank of \(A \otimes_R M\) at \(\mathfrak q\) equals the fibre rank
  of \(M\) at \(\mathfrak p\).  No flatness or finiteness hypotheses are
  needed.
\end{theorem}

\begin{proof}
  \leanok
  The residue-field map \(R \to \kappa(\mathfrak p) \to \kappa(\mathfrak q)\)
  makes \(\kappa(\mathfrak q)\) a \(\kappa(\mathfrak p)\)-algebra.  Cancelling
  base changes gives \(\kappa(\mathfrak p)\)-linear, then
  \(\kappa(\mathfrak q)\)-linear, isomorphisms
  \[
    \kappa(\mathfrak q) \otimes_A (A \otimes_R M)
      \;\cong\; \kappa(\mathfrak q) \otimes_R M
      \;\cong\; \kappa(\mathfrak q) \otimes_{\kappa(\mathfrak p)}
        \bigl(\kappa(\mathfrak p) \otimes_R M\bigr).
  \]
  Extension of scalars along a field extension preserves dimension, so
  \(\dim_{\kappa(\mathfrak q)}
    \kappa(\mathfrak q) \otimes_{\kappa(\mathfrak p)}
    (\kappa(\mathfrak p) \otimes_R M)
    = \dim_{\kappa(\mathfrak p)}(\kappa(\mathfrak p) \otimes_R M)\),
  which is the claim.
\end{proof}

\begin{lemma}[Fibre rank is a linear invariant]
  \label{lem:fiber_rank_congr}
  \dcref{custom}
  \lean{Ideal.fiberRank_congr}
  \leanok
  \uses{def:fiber_rank}
  If \(\sigma : M \cong N\) is an isomorphism of \(R\)-modules, then \(M\) and
  \(N\) have the same fibre rank at every prime \(\mathfrak p\).
\end{lemma}

\begin{proof}
  \leanok
  Base change \(\sigma\) along \(R \to \kappa(\mathfrak p)\) to obtain a
  \(\kappa(\mathfrak p)\)-linear isomorphism
  \(\kappa(\mathfrak p) \otimes_R M \cong \kappa(\mathfrak p) \otimes_R N\);
  isomorphic vector spaces have equal dimension.
\end{proof}

\begin{corollary}[Fibre rank of a localized module]
  \label{cor:fiber_rank_localized}
  \dcref{custom}
  \lean{Ideal.fiberRank_of_isLocalizedModule}
  \leanok
  \uses{thm:fiber_rank_base_change, lem:fiber_rank_congr}
  Let \(R' = S^{-1}R\) be a localization, \(N = S^{-1}M\) the corresponding
  localized module, and \(\mathfrak q\) a prime of \(R'\).  Then the fibre
  rank of \(N\) at \(\mathfrak q\) equals the fibre rank of \(M\) at the
  contraction \(\mathfrak q \cap R\).
\end{corollary}

\begin{proof}
  \leanok
  Localization is a base change: the canonical map \(M \to N\) induces an
  \(R'\)-linear isomorphism \(R' \otimes_R M \cong N\).  By
  \cref{lem:fiber_rank_congr} the fibre rank of \(N\) at \(\mathfrak q\)
  equals that of \(R' \otimes_R M\), which by \cref{thm:fiber_rank_base_change}
  equals the fibre rank of \(M\) at \(\mathfrak q \cap R\).
\end{proof}

\subsection{Presentation charts and the point rank}

\begin{definition}[Presentation chart]
  \label{def:presentation_chart}
  \lean{AlgebraicGeometry.Scheme.Modules.IsPresentationChart}
  \leanok
  \uses{def:matrix_presentation}
  \dcref{custom}
  An affine open \(V \subseteq X\) is an \emph{\(e\)-presentation chart} for
  \(\mathcal G\) if the section module \(\Gamma(\mathcal G, V)\) admits a
  matrix presentation by \(e\) generators over \(\Gamma(X, V)\), that is, a
  right-exact sequence
  \(\OO_V^{\oplus m} \xrightarrow{\psi} \OO_V^{\oplus e}
    \to \mathcal G|_V \to 0\).
\end{definition}

\begin{definition}[Restriction of a presentation to a basic open]
  \label{def:presentation_basic_open}
  \dcref{custom}
  \lean{AlgebraicGeometry.Scheme.Modules.MatrixPresentationBasicOpen}
  \leanok
  \uses{def:matrix_presentation_base_change, def:presentation_chart}
  Let \(P = (\psi, \pi)\) be a matrix presentation of
  \(\Gamma(\mathcal G, V)\) by \(e\) generators over an affine open \(V\), and
  \(f \in \Gamma(X, V)\).  Since \(\mathcal G\) is quasi-coherent, restriction
  exhibits \(\Gamma(\mathcal G, D(f))\) as the localization
  \(\Gamma(\mathcal G, V) \otimes_{\Gamma(X, V)} \Gamma(X, D(f))\).  Base
  change of \(P\) along \(\Gamma(X, V) \to \Gamma(X, D(f))\)
  (\cref{def:matrix_presentation_base_change}), transported along this
  localization isomorphism, is a matrix presentation of
  \(\Gamma(\mathcal G, D(f))\) by \(e\) generators, the \emph{restriction of
  \(P\) to \(D(f)\)}; its relation matrix is the entrywise restriction
  \(\psi|_{D(f)}\).
\end{definition}

\begin{lemma}[Entry ideal of a restricted presentation]
  \label{lem:entry_ideal_basic_open}
  \dcref{custom}
  \lean{AlgebraicGeometry.Scheme.Modules.entryIdeal_matrixPresentationBasicOpen}
  \leanok
  \uses{def:presentation_basic_open, lem:entry_ideal_base_change}
  With the notation of \cref{def:presentation_basic_open}, the entry ideal of
  the restricted presentation is the image of \(I(\psi)\) under restriction:
  \(I(\psi|_{D(f)}) = I(\psi) \cdot \Gamma(X, D(f))\).
\end{lemma}

\begin{proof}
  \leanok
  The restricted presentation is a base change of \(P\), so its entry ideal is
  the image ideal by \cref{lem:entry_ideal_base_change}.
\end{proof}

\begin{lemma}[Presentation charts pass to basic opens]
  \label{lem:presentation_chart_basic_open}
  \dcref{custom}
  \lean{AlgebraicGeometry.Scheme.Modules.IsPresentationChart.basicOpen}
  \leanok
  \uses{def:presentation_chart, def:presentation_basic_open}
  If \(V\) is an \(e\)-presentation chart for \(\mathcal G\) and
  \(f \in \Gamma(X, V)\), then the basic open \(D(f) \subseteq V\) is again an
  \(e\)-presentation chart.
\end{lemma}

\begin{proof}
  \leanok
  Restrict a presentation of \(\Gamma(\mathcal G, V)\) to \(D(f)\) by
  \cref{def:presentation_basic_open}.
\end{proof}

\begin{definition}[Chart fibre rank]
  \label{def:chart_fiber_rank}
  \dcref{custom}
  \lean{AlgebraicGeometry.Scheme.Modules.chartFiberRank}
  \leanok
  \uses{def:fiber_rank}
  For an affine open chart \(V\) and a point \(x \in V\), the \emph{chart
  fibre rank} of \(\mathcal G\) at \(x\) computed on \(V\) is the fibre rank
  of the section module \(\Gamma(\mathcal G, V)\) at the prime
  \(\mathfrak p_x \subseteq \Gamma(X, V)\) corresponding to \(x\).
\end{definition}

\begin{lemma}[Chart fibre rank under localization]
  \label{lem:chart_fiber_rank_basic_open}
  \dcref{custom}
  \lean{AlgebraicGeometry.Scheme.Modules.chartFiberRank_basicOpen}
  \leanok
  \uses{def:chart_fiber_rank, cor:fiber_rank_localized}
  The chart fibre rank is unchanged by passing to a basic open of the chart:
  for \(f \in \Gamma(X, V)\) and \(x \in D(f)\), the chart fibre rank computed
  on \(V\) equals that computed on \(D(f)\).
\end{lemma}

\begin{proof}
  \leanok
  \(\Gamma(\mathcal G, D(f))\) is the localization of
  \(\Gamma(\mathcal G, V)\) at \(f\), and the prime of \(x\) in
  \(\Gamma(X, D(f))\) contracts to the prime of \(x\) in \(\Gamma(X, V)\).
  Apply \cref{cor:fiber_rank_localized}.
\end{proof}

\begin{theorem}[Chart-independence of the fibre rank]
  \label{thm:chart_fiber_rank_eq}
  \dcref{custom}
  \lean{AlgebraicGeometry.Scheme.Modules.chartFiberRank_eq}
  \leanok
  \uses{def:chart_fiber_rank, lem:chart_fiber_rank_basic_open}
  Any two affine charts \(V, W\) containing a point \(x\) give the same chart
  fibre rank of \(\mathcal G\) at \(x\).
\end{theorem}

\begin{proof}
  \leanok
  Choose a basic open of \(V\) and a basic open of \(W\) that coincide as an
  open \(D \ni x\) contained in \(V \cap W\).  By
  \cref{lem:chart_fiber_rank_basic_open} the chart fibre rank computed on
  \(V\) and the one computed on \(W\) both equal the chart fibre rank computed
  on \(D\), hence agree.
\end{proof}

\begin{definition}[Point rank]
  \label{def:point_rank}
  \dcref{custom}
  \lean{AlgebraicGeometry.Scheme.Modules.pointRank}
  \leanok
  \uses{def:chart_fiber_rank, thm:chart_fiber_rank_eq}
  The \emph{point rank} of \(\mathcal G\) at a point \(x \in X\) is the
  dimension of the fibre of \(\mathcal G\) at \(x\), computed as the chart
  fibre rank on any affine chart through \(x\); it is well defined by
  \cref{thm:chart_fiber_rank_eq}.
\end{definition}

\begin{lemma}[Point rank on a chart]
  \label{lem:point_rank_eq_chart_fiber_rank}
  \dcref{custom}
  \lean{AlgebraicGeometry.Scheme.Modules.pointRank_eq_chartFiberRank}
  \leanok
  \uses{def:point_rank, thm:chart_fiber_rank_eq}
  For any affine chart \(V \ni x\), the point rank of \(\mathcal G\) at \(x\)
  equals the chart fibre rank computed on \(V\).
\end{lemma}

\begin{proof}
  \leanok
  Both are chart fibre ranks at \(x\); apply \cref{thm:chart_fiber_rank_eq}.
\end{proof}

\subsection{The strata ideal}

\begin{lemma}[Locality of entry-ideal membership]
  \label{lem:mem_entry_ideal_locally}
  \dcref{custom}
  \lean{AlgebraicGeometry.Scheme.Modules.mem_entryIdeal_of_locally}
  \leanok
  \uses{def:entry_ideal, lem:entry_ideal_basic_open}
  Let \(P = (\psi, \pi)\) be a matrix presentation of
  \(\Gamma(\mathcal G, V)\) by \(e\) generators over an affine open \(V\), and
  \(r \in \Gamma(X, V)\).  If every point of \(V\) has a basic open
  \(D(g) \ni x\) on which the restriction of \(r\) lies in the entry ideal of
  the restricted presentation \(\psi|_{D(g)}\), then \(r \in I(\psi)\).
\end{lemma}

\begin{proof}
  \leanok
  Choose for each point such a \(g\); the basic opens \(D(g)\) cover \(V\), so
  the corresponding \(g\) generate the unit ideal of \(\Gamma(X, V)\).  By
  \cref{lem:entry_ideal_basic_open}, membership of \(r|_{D(g)}\) in
  \(I(\psi|_{D(g)}) = I(\psi) \cdot \Gamma(X, D(g))\) means a power
  \(g^{n} r\) lies in \(I(\psi)\).  A section that, after multiplication by a
  power of each member of a family generating the unit ideal, lands in a
  submodule, already lies in that submodule; hence \(r \in I(\psi)\).
\end{proof}

\begin{lemma}[Chart-to-chart transfer of the entry-ideal condition]
  \label{lem:entry_ideal_chart_transfer}
  \dcref{custom}
  \lean{AlgebraicGeometry.Scheme.Modules.exists_basicOpen_mem_entryIdeal_of_chart}
  \leanok
  \uses{thm:entry_ideal_presentation_independent, lem:entry_ideal_basic_open}
  Let \(r \in \Gamma(X, U)\), let \(V, W \subseteq U\) be affine opens with
  \(e\)-generator presentations \(P_V, P_W\), and suppose the restriction of
  \(r\) to \(W\) lies in \(I(\psi_W)\).  Then for every point
  \(x \in V \cap W\) there is a basic open \(D(g) \subseteq V\) with
  \(x \in D(g)\) on which the restriction of \(r\) lies in
  \(I(\psi_V|_{D(g)})\).
\end{lemma}

\begin{proof}
  \leanok
  Choose a common basic open \(D \ni x\) contained in \(V \cap W\), realised
  as a basic open \(D(g)\) of \(V\) and a basic open \(D(f)\) of \(W\).  On
  \(D\) both \(\psi_V|_{D(g)}\) and \(\psi_W|_{D(f)}\) present
  \(\Gamma(\mathcal G, D)\) by \(e\) generators, so by
  \cref{thm:entry_ideal_presentation_independent} their entry ideals coincide.
  By \cref{lem:entry_ideal_basic_open} the restriction of \(r|_W\) to \(D\)
  lies in \(I(\psi_W|_{D(f)}) = I(\psi_V|_{D(g)})\).
\end{proof}

\begin{definition}[Strata ideal]
  \label{def:strata_ideal}
  \lean{AlgebraicGeometry.Scheme.Modules.strataIdeal}
  \leanok
  \uses{def:entry_ideal, def:presentation_chart}
  \dcref{custom}
  The rank-\(e\) \emph{strata ideal} of \(\mathcal G\) on an affine open \(U\)
  is the ideal of \(\Gamma(X, U)\)
  \[
    I_e(U) \;=\; \bigcap_{V \le U} \; \bigcap_{P}
      \; \bigl(\text{restriction } \Gamma(X, U) \to \Gamma(X, V)\bigr)^{-1}
        \bigl(I(\psi_P)\bigr),
  \]
  the intersection over all affine opens \(V \subseteq U\) and all matrix
  presentations \(P = (\psi_P, \pi_P)\) of \(\Gamma(\mathcal G, V)\) by \(e\)
  generators, of the preimages of the entry ideals under restriction.
\end{definition}

\begin{theorem}[Strata ideal on a chart is the entry ideal]
  \label{thm:strata_ideal_eq_entry_ideal}
  \dcref{custom}
  \lean{AlgebraicGeometry.Scheme.Modules.strataIdeal_eq_entryIdeal}
  \leanok
  \uses{def:strata_ideal, lem:mem_entry_ideal_locally, lem:entry_ideal_chart_transfer}
  If \(V\) is an \(e\)-presentation chart with presentation \(P\), then the
  strata ideal on \(V\) is the entry ideal of \(P\): \(I_e(V) = I(\psi_P)\).
\end{theorem}

\begin{proof}
  \leanok
  For \(\subseteq\): the pair \((V, P)\) itself occurs in the defining
  intersection, and restriction along \(V \le V\) is the identity, so every
  element of \(I_e(V)\) lies in \(I(\psi_P)\).  For \(\supseteq\): let
  \(r \in I(\psi_P)\) and let \((W, Q)\) be any presentation chart
  \(W \subseteq V\).  By \cref{lem:mem_entry_ideal_locally} it suffices to
  show that each point of \(W\) has a basic open on which \(r|_W\) lands in the
  entry ideal of the restricted \(Q\); this is
  \cref{lem:entry_ideal_chart_transfer} applied to the two charts \(W\) and
  \(V\) (using \(r|_V = r \in I(\psi_P)\)).  Hence \(r|_W \in I(\psi_Q)\) for
  every chart, i.e.\ \(r \in I_e(V)\).
\end{proof}

\begin{lemma}[Strata membership from a covering family of charts]
  \label{lem:mem_strata_ideal_of_charts}
  \dcref{custom}
  \lean{AlgebraicGeometry.Scheme.Modules.mem_strataIdeal_of_charts}
  \leanok
  \uses{def:strata_ideal, lem:mem_entry_ideal_locally, lem:entry_ideal_chart_transfer}
  Let \(r \in \Gamma(X, U)\) and let \((V_i)\) be a family of affine
  presentation charts \(V_i \subseteq U\), with presentations \(P_i\), whose
  underlying opens cover \(U\).  If the restriction of \(r\) to each \(V_i\)
  lies in \(I(\psi_{P_i})\), then \(r \in I_e(U)\).
\end{lemma}

\begin{proof}
  \leanok
  Let \((W, Q)\) be any presentation chart \(W \subseteq U\).  By
  \cref{lem:mem_entry_ideal_locally} it suffices to check local membership at
  each \(x \in W\); choosing an index \(i\) with \(x \in V_i\),
  \cref{lem:entry_ideal_chart_transfer} applied to \(W\) and \(V_i\) (using
  \(r|_{V_i} \in I(\psi_{P_i})\)) supplies a basic open of \(W\) on which
  \(r\) lands in the entry ideal of the restricted \(Q\).  Thus
  \(r|_W \in I(\psi_Q)\) for every chart, i.e.\ \(r \in I_e(U)\).
\end{proof}

\begin{definition}[Charts-cover hypothesis]
  \label{def:charts_cover}
  \dcref{custom}
  \lean{AlgebraicGeometry.Scheme.Modules.ChartsCover}
  \leanok
  \uses{def:presentation_chart}
  The module \(\mathcal G\) \emph{admits \(e\)-presentation charts} if every
  point of \(X\) lies in an \(e\)-presentation chart.  This is the standing
  hypothesis of the rank-\(e\) stratum construction; it holds on the union of
  all \(e\)-presentation charts.
\end{definition}

\begin{lemma}[Finite chart covers]
  \label{lem:charts_cover_finite}
  \dcref{custom}
  \lean{AlgebraicGeometry.Scheme.Modules.ChartsCover.exists_finite_charts}
  \leanok
  \uses{def:charts_cover, lem:presentation_chart_basic_open}
  If \(\mathcal G\) admits \(e\)-presentation charts, then every affine open
  \(U \subseteq X\) is covered by finitely many \(e\)-presentation charts
  contained in \(U\).
\end{lemma}

\begin{proof}
  \leanok
  Each point of \(U\) lies in some \(e\)-presentation chart; refining to a
  basic open contained in \(U\) keeps it a presentation chart
  (\cref{lem:presentation_chart_basic_open}).  These charts cover \(U\), and
  \(U\) is quasi-compact, being affine, so finitely many of them suffice.
\end{proof}

\begin{theorem}[The strata ideal is quasi-coherent]
  \label{thm:map_strata_ideal_basic_open}
  \lean{AlgebraicGeometry.Scheme.Modules.map_strataIdeal_basicOpen}
  \leanok
  \uses{def:strata_ideal, lem:charts_cover_finite, lem:mem_strata_ideal_of_charts,
    lem:entry_ideal_basic_open, thm:strata_ideal_eq_entry_ideal}
  \dcref{custom}
  Suppose \(\mathcal G\) admits \(e\)-presentation charts.  For an affine open
  \(U\) and \(f \in \Gamma(X, U)\), restriction along
  \(\Gamma(X, U) \to \Gamma(X, D(f))\) carries the strata ideal to the strata
  ideal: the image of \(I_e(U)\) generates \(I_e(D(f))\).  Consequently the
  assignment \(U \mapsto I_e(U)\) is quasi-coherent.
\end{theorem}

\begin{proof}
  \leanok
  For \(\subseteq\): any presentation chart \(W \subseteq D(f)\) is also a
  chart in \(U\), so a section of \(I_e(U)\) restricts into \(I(\psi_W)\);
  hence its restriction to \(D(f)\) lies in \(I_e(D(f))\).

  For \(\supseteq\): let \(r' \in I_e(D(f))\) and write
  \(r' = y / f^{n}\) with \(y \in \Gamma(X, U)\).  By
  \cref{lem:charts_cover_finite} choose finitely many presentation charts
  \(V_1, \dots, V_N\) covering \(U\), with presentations \(P_i\).  On each
  \(V_i\), the basic open \(D(f|_{V_i}) \subseteq D(f)\) is a chart, and the
  strata condition on \(r'\) there says \(y / f^{n}\) restricts into
  \(I(\psi_{P_i}|_{D(f|_{V_i})})\), the image ideal
  (\cref{lem:entry_ideal_basic_open}); clearing the denominator yields an
  exponent \(m_i\) with \(f^{m_i} y|_{V_i} \in I(\psi_{P_i})\).  Setting
  \(m = \max_i m_i\), multiplying up gives
  \(f^{m} y|_{V_i} \in I(\psi_{P_i})\) for all \(i\), so
  \(f^{m} y \in I_e(U)\) by \cref{lem:mem_strata_ideal_of_charts}.  Therefore
  \(r' = (f^{m} y) / f^{m + n}\) lies in the image of \(I_e(U)\) under
  restriction.  (\cref{thm:strata_ideal_eq_entry_ideal} identifies the strata
  ideal with the entry ideal on each chart, licensing the denominator
  clearing.)
\end{proof}

\subsection{The rank stratum}

\begin{definition}[Strata ideal-sheaf datum]
  \label{def:strata_data}
  \dcref{custom}
  \lean{AlgebraicGeometry.Scheme.Modules.strataData}
  \leanok
  \uses{def:strata_ideal, thm:map_strata_ideal_basic_open}
  Suppose \(\mathcal G\) admits \(e\)-presentation charts.  The rank-\(e\)
  \emph{strata ideal-sheaf datum} on \(X\) is the assignment
  \(U \mapsto I_e(U)\) of \cref{def:strata_ideal}, which is a quasi-coherent
  ideal sheaf by \cref{thm:map_strata_ideal_basic_open}.
\end{definition}

\begin{definition}[Rank-\(e\) stratum]
  \label{def:stratum}
  \lean{AlgebraicGeometry.Scheme.Modules.stratum}
  \leanok
  \uses{def:strata_data}
  \dcref{custom}
  The rank-\(e\) \emph{stratum} of \(\mathcal G\) is the closed subscheme of
  \(X\) cut out by the strata ideal-sheaf datum of \cref{def:strata_data}.
\end{definition}

\begin{definition}[Closed immersion of the stratum]
  \label{def:stratum_immersion}
  \dcref{custom}
  \lean{AlgebraicGeometry.Scheme.Modules.stratumι}
  \leanok
  \uses{def:stratum}
  The canonical morphism from the rank-\(e\) stratum to \(X\) is a closed
  immersion.
\end{definition}

\begin{theorem}[Support of the rank-\(e\) stratum]
  \label{thm:mem_range_stratum}
  \lean{AlgebraicGeometry.Scheme.Modules.mem_range_stratumι_iff}
  \leanok
  \uses{def:stratum_immersion, thm:strata_ideal_eq_entry_ideal, thm:entry_ideal_prime_iff,
    def:point_rank, lem:point_rank_eq_chart_fiber_rank}
  \dcref{custom}
  A point \(x \in X\) lies in the image of the closed immersion of the
  rank-\(e\) stratum if and only if the point rank of \(\mathcal G\) at \(x\)
  equals \(e\); that is, the support of the stratum is exactly the locus where
  the fibre of \(\mathcal G\) has dimension exactly \(e\).
\end{theorem}

\begin{proof}
  \leanok
  Choose an \(e\)-presentation chart \(V \ni x\) with presentation \(P\).  The
  image of the closed immersion is the support of the ideal sheaf, which meets
  \(V\) in the zero locus of \(I_e(V)\); and \(I_e(V) = I(\psi_P)\) by
  \cref{thm:strata_ideal_eq_entry_ideal}.  A point of an affine open lies in
  the zero locus of an ideal of sections iff its prime contains that ideal, so
  \(x\) lies in the stratum iff \(I(\psi_P) \subseteq \mathfrak p_x\).  By
  \cref{thm:entry_ideal_prime_iff} this holds iff the fibre rank of
  \(\Gamma(\mathcal G, V)\) at \(\mathfrak p_x\) --- the chart fibre rank at
  \(x\), equal to the point rank by
  \cref{lem:point_rank_eq_chart_fiber_rank} --- equals \(e\).
\end{proof}

\section{Representability of the rank strata}
\label{sec:flatstrat_representability}

The rank-\(e\) stratum of \cref{sec:flatstrat_rank_stratum} is a closed
subscheme with the correct support.  This section establishes that it also
has the correct functor of points --- Nitsure's local universal property:
the pullback of \(\mathcal G\) along a morphism is locally free of rank
\(e\) exactly when the pulled-back relation matrices vanish, exactly when
the morphism factors through the stratum --- and assembles the strata over
the chart loci into the locally closed pieces of the flattening
stratification of a finitely presented module on a locally noetherian
scheme.  Throughout, ``flat over \(T\)'' for an \(\OO_T\)-module means flat
over \(T\) via the identity morphism of \(T\), in the affine-local sense of
\cref{def:coherent_sheaf_flat}.

\subsection{Pullback presentations and flatness of sections}

\begin{lemma}[Finitely presented modules admit matrix presentations]
  \label{lem:exists_matrix_presentation}
  \dcref{custom}
  \lean{Module.FinitePresentation.exists_matrixPresentation}
  \leanok
  \uses{def:matrix_presentation}
  Every finitely presented module \(M\) over a commutative ring \(R\) admits
  a matrix presentation: there are natural numbers \(e, m\) and a matrix
  presentation of \(M\) by \(e\) generators and \(m\) relations.
\end{lemma}

\begin{proof}
  \leanok
  By finite presentation there is a surjection \(\pi : R^{\oplus e} \to M\)
  whose kernel \(K\) is finitely generated; choose generators
  \(t_1, \dots, t_m\) of \(K \subseteq R^{\oplus e}\) and let \(\psi\) be
  the \(e \times m\) matrix with columns \(t_j\).  The image of
  \(\psi : R^{\oplus m} \to R^{\oplus e}\) is the submodule generated by
  the columns, namely \(K = \ker \pi\), so
  \(R^{\oplus m} \xrightarrow{\psi} R^{\oplus e} \xrightarrow{\pi} M \to 0\)
  is exact.
\end{proof}

\begin{lemma}[Pullback of a presentation to section modules]
  \label{lem:presentation_pullback_sections}
  \lean{AlgebraicGeometry.Scheme.Modules.exists_matrixPresentation_pullback_sections}
  \leanok
  \uses{def:matrix_presentation}
  \dcref{custom}
  Let \(g : Y \to X\) be a morphism of schemes, \(\mathcal G\) a
  quasi-coherent \(\OO_X\)-module, \(V \subseteq X\) an affine open,
  \(W \subseteq g^{-1}(V)\) an affine open, and \(P = (\psi, \pi)\) a matrix
  presentation of \(\Gamma(\mathcal G, V)\) by \(e\) generators and \(m\)
  relations over \(\Gamma(X, V)\).  Then \(\Gamma(g^*\mathcal G, W)\) admits
  a matrix presentation by \(e\) generators and \(m\) relations over
  \(\Gamma(Y, W)\) whose relation matrix is the entrywise image
  \(g^\sharp(\psi)\) of \(\psi\) under the induced ring map
  \(g^\sharp : \Gamma(X, V) \to \Gamma(Y, W)\).
\end{lemma}

\begin{proof}
  \leanok
  \uses{def:pullback_app_isoTensor_sigma, def:matrix_presentation_base_change}
  For the affine pair \((V, W)\) the sections of the pullback module satisfy
  the base-change formula
  \(\Gamma(g^*\mathcal G, W) \cong
    \Gamma(Y, W) \otimes_{\Gamma(X, V)} \Gamma(\mathcal G, V)\)
  as \(\Gamma(Y, W)\)-modules (\cref{def:pullback_app_isoTensor_sigma}).
  Base change of \(P\) along \(g^\sharp\)
  (\cref{def:matrix_presentation_base_change}) is a matrix presentation of
  the right-hand side --- the pulled-back sequence stays exact by
  right-exactness of the tensor product --- with relation matrix
  \(g^\sharp(\psi)\), and transporting it along the base-change isomorphism
  changes neither the generator count nor the relation matrix.
\end{proof}

\begin{lemma}[Flat sections from flatness over the identity]
  \label{lem:flat_sections_of_sheaf_flat}
  \dcref{custom}
  \lean{AlgebraicGeometry.Scheme.Modules.flat_sections_of_coherentSheafFlat_id}
  \leanok
  \uses{def:coherent_sheaf_flat}
  Let \(\mathcal G\) be an \(\OO_T\)-module that is flat over \(T\).  Then
  for every affine open \(W \subseteq T\) the section module
  \(\Gamma(\mathcal G, W)\) is flat over \(\Gamma(T, W)\).
\end{lemma}

\begin{proof}
  \leanok
  Specialize the affine-local form of flatness over the identity to the pair
  of affine opens \(U := W\) and \(W \subseteq \mathrm{id}^{-1}(U) = U\):
  the module \(\Gamma(\mathcal G, W)\) is flat over \(\Gamma(T, W)\) acting
  through the ring map on sections induced by the identity morphism, which
  is the identity of \(\Gamma(T, W)\), so the action is the canonical one.
\end{proof}

\begin{lemma}[Flatness over the identity from an affine cover]
  \label{lem:sheaf_flat_of_flat_charts}
  \dcref{custom}
  \lean{AlgebraicGeometry.Scheme.Modules.coherentSheafFlat_id_of_charts}
  \leanok
  \uses{def:coherent_sheaf_flat}
  Let \(\mathcal G\) be a quasi-coherent \(\OO_T\)-module and \((W_j)_j\) an
  affine open cover of \(T\) such that each \(\Gamma(\mathcal G, W_j)\) is
  flat over \(\Gamma(T, W_j)\).  Then \(\mathcal G\) is flat over \(T\).
\end{lemma}

\begin{proof}
  \leanok
  \uses{lem:flat_affine_local}
  Apply the affine-locality of flatness over the base
  (\cref{lem:flat_affine_local}) to the identity morphism of \(T\) with the
  charts \(U_j := W_j\), \(W_j \subseteq \mathrm{id}^{-1}(U_j)\): flatness
  of the section modules on the covering charts gives flatness of
  \(\Gamma(\mathcal G, W)\) over \(\Gamma(T, U)\) for every pair of affine
  opens \(W \subseteq U\), which is flatness over the identity.
\end{proof}

\begin{lemma}[Flatness over the base is stable under pullback]
  \label{lem:sheaf_flat_pullback}
  \dcref{custom}
  \lean{AlgebraicGeometry.Scheme.Modules.coherentSheafFlat_id_pullback}
  \leanok
  \uses{def:coherent_sheaf_flat, lem:pullback_isQuasicoherent_hom}
  Let \(\mathcal G\) be a quasi-coherent \(\OO_T\)-module flat over \(T\),
  and \(g : T' \to T\) a morphism of schemes.  Then \(g^*\mathcal G\) is
  flat over \(T'\).
\end{lemma}

\begin{proof}
  \leanok
  \uses{lem:flat_sections_of_sheaf_flat, def:pullback_app_isoTensor_sigma,
    lem:sheaf_flat_of_flat_charts}
  Cover \(T'\) by affine opens \(W\), each contained in the preimage of some
  affine open \(V \subseteq T\).  For such a pair the section module
  \(\Gamma(g^*\mathcal G, W) \cong
    \Gamma(T', W) \otimes_{\Gamma(T, V)} \Gamma(\mathcal G, V)\)
  (\cref{def:pullback_app_isoTensor_sigma}) is a base change of the flat
  module \(\Gamma(\mathcal G, V)\) (\cref{lem:flat_sections_of_sheaf_flat}),
  hence flat over \(\Gamma(T', W)\).  The pullback \(g^*\mathcal G\) is
  quasi-coherent, so \cref{lem:sheaf_flat_of_flat_charts} upgrades the cover
  of flat section modules to flatness over \(T'\).
\end{proof}

\subsection{The rank dictionary}

\begin{theorem}[Stalk rank of an affine section module]
  \label{thm:stalk_rank_eq_point_rank}
  \dcref{custom}
  \lean{AlgebraicGeometry.Scheme.Modules.rankAtStalk_sections_eq_pointRank}
  \leanok
  \uses{def:point_rank}
  Let \(\mathcal G\) be a quasi-coherent \(\OO_T\)-module and
  \(W \subseteq T\) an affine open such that \(M := \Gamma(\mathcal G, W)\)
  is flat and finite over \(R := \Gamma(T, W)\).  Let \(\mathfrak p\) be a
  prime of \(R\) and \(x \in W\) the corresponding point under the canonical
  identification of \(W\) with the prime spectrum of \(R\).  Then the
  localization \(M_{\mathfrak p}\) is a free \(R_{\mathfrak p}\)-module
  whose rank equals the point rank of \(\mathcal G\) at \(x\).
\end{theorem}

\begin{proof}
  \leanok
  \uses{def:fiber_rank, def:chart_fiber_rank,
    lem:point_rank_eq_chart_fiber_rank}
  \(M_{\mathfrak p}\) is finite and flat over the local ring
  \(R_{\mathfrak p}\), hence free of some rank \(r\) (a finitely generated
  flat module over a local ring is free).  Reducing modulo the maximal
  ideal,
  \(M \otimes_R \kappa(\mathfrak p)
    \cong M_{\mathfrak p} \otimes_{R_{\mathfrak p}} \kappa(\mathfrak p)\)
  has dimension \(r\) over \(\kappa(\mathfrak p)\), so \(r\) is the fiber
  rank of \(M\) at \(\mathfrak p\) (\cref{def:fiber_rank}).  The prime of
  \(R\) attached to the point \(x\) is \(\mathfrak p\) itself, so this fiber
  rank is the chart fibre rank of \(\mathcal G\) at \(x\) computed on \(W\)
  (\cref{def:chart_fiber_rank}), which equals the point rank by
  \cref{lem:point_rank_eq_chart_fiber_rank}.
\end{proof}

\begin{theorem}[Point rank of a pullback]
  \label{thm:point_rank_pullback}
  \dcref{custom}
  \lean{AlgebraicGeometry.Scheme.Modules.pointRank_pullback}
  \leanok
  \uses{def:point_rank, lem:pullback_isQuasicoherent_hom}
  Let \(g : Y \to X\) be a morphism of schemes and \(\mathcal G\) a
  quasi-coherent \(\OO_X\)-module.  For every \(y \in Y\), the point rank of
  \(g^*\mathcal G\) at \(y\) equals the point rank of \(\mathcal G\) at
  \(g(y)\).  No flatness or finiteness hypotheses are needed.
\end{theorem}

\begin{proof}
  \leanok
  \uses{def:pullback_app_isoTensor_sigma, thm:fiber_rank_base_change,
    lem:fiber_rank_congr, lem:point_rank_eq_chart_fiber_rank}
  Choose an affine open \(V \ni g(y)\) and an affine open \(W \ni y\) with
  \(W \subseteq g^{-1}(V)\).  By the base-change formula for sections
  (\cref{def:pullback_app_isoTensor_sigma}),
  \(\Gamma(g^*\mathcal G, W) \cong
    \Gamma(Y, W) \otimes_{\Gamma(X, V)} \Gamma(\mathcal G, V)\).  The prime
  \(\mathfrak q \subseteq \Gamma(Y, W)\) of \(y\) contracts along the
  induced ring map to the prime \(\mathfrak p \subseteq \Gamma(X, V)\) of
  \(g(y)\).  By \cref{lem:fiber_rank_congr} and base-change invariance of
  the fiber rank (\cref{thm:fiber_rank_base_change}), the fiber rank of
  \(\Gamma(g^*\mathcal G, W)\) at \(\mathfrak q\) equals the fiber rank of
  \(\Gamma(\mathcal G, V)\) at \(\mathfrak p\).  These are the chart fibre
  ranks of \(g^*\mathcal G\) at \(y\) on \(W\) and of \(\mathcal G\) at
  \(g(y)\) on \(V\), i.e.\ the two point ranks
  (\cref{lem:point_rank_eq_chart_fiber_rank}).
\end{proof}

\subsection{Factorization through the stratum}

\begin{theorem}[The strata ideal dies under a flat constant-rank pullback]
  \label{thm:strata_ideal_le_ker}
  \lean{AlgebraicGeometry.Scheme.Modules.strataData_le_ker}
  \leanok
  \uses{def:strata_data, def:charts_cover, def:coherent_sheaf_flat,
    def:point_rank, lem:pullback_isQuasicoherent_hom}
  \dcref{custom}
  Let \(\mathcal G\) be a quasi-coherent \(\OO_X\)-module admitting
  \(e\)-presentation charts, and let \(q : T \to X\) be a morphism such that
  \(q^*\mathcal G\) is flat over \(T\) and has point rank \(e\) at every
  point of \(T\).  Then the strata ideal sheaf is contained in the kernel of
  \(q\): for every affine open \(U \subseteq X\), every \(r \in I_e(U)\)
  pulls back to \(0 \in \Gamma(T, q^{-1}(U))\).
\end{theorem}

\begin{proof}
  \leanok
  \uses{lem:charts_cover_finite, lem:presentation_pullback_sections,
    lem:flat_sections_of_sheaf_flat, thm:stalk_rank_eq_point_rank,
    thm:entry_ideal_flat_vanish, lem:entry_ideal_vanish, def:strata_ideal,
    lem:appLE_base_res_apply}
  A section of a sheaf vanishes if it vanishes on a cover, so it suffices to
  find, around each point \(t \in q^{-1}(U)\), an open on which
  \(q^\sharp(r)\) restricts to \(0\).  Choose finitely many
  \(e\)-presentation charts \(V_1, \dots, V_N \subseteq U\) covering \(U\),
  with presentations \(P_i = (\psi_i, \pi_i)\)
  (\cref{lem:charts_cover_finite}).  Given \(t\), pick \(i\) with
  \(q(t) \in V_i\), and an affine open \(W \ni t\) with
  \(W \subseteq q^{-1}(V_i)\).

  By \cref{lem:presentation_pullback_sections} the module
  \(\Gamma(q^*\mathcal G, W)\) has a matrix presentation by \(e\) generators
  with relation matrix \(q^\sharp(\psi_i)\); in particular it is finite over
  \(\Gamma(T, W)\), and it is flat by
  \cref{lem:flat_sections_of_sheaf_flat}.  At every prime of
  \(\Gamma(T, W)\) its localization is free of rank equal to the point rank
  of \(q^*\mathcal G\) at the corresponding point
  (\cref{thm:stalk_rank_eq_point_rank}), which is \(e\) by hypothesis.  By
  \cref{thm:entry_ideal_flat_vanish} the relation matrix vanishes:
  \(q^\sharp(\psi_i) = 0\).  Hence the entry ideal \(I(\psi_i)\) is
  contained in the kernel of
  \(q^\sharp : \Gamma(X, V_i) \to \Gamma(T, W)\)
  (\cref{lem:entry_ideal_vanish}).  The strata-ideal condition
  (\cref{def:strata_ideal}) gives \(r|_{V_i} \in I(\psi_i)\), and the maps
  induced by \(q\) on sections commute with restriction
  (\cref{lem:appLE_base_res_apply}), so
  \(q^\sharp(r)|_W = q^\sharp(r|_{V_i}) = 0\), as required.
\end{proof}

\begin{theorem}[Unique factorization through the rank-\(e\) stratum]
  \label{thm:stratum_lift_unique}
  \lean{AlgebraicGeometry.Scheme.Modules.existsUnique_stratumLift}
  \leanok
  \uses{def:stratum, def:stratum_immersion, def:charts_cover,
    def:coherent_sheaf_flat, def:point_rank, lem:pullback_isQuasicoherent_hom}
  \dcref{custom}
  Let \(\mathcal G\) be a quasi-coherent \(\OO_X\)-module admitting
  \(e\)-presentation charts, and let \(q : T \to X\) be a morphism such that
  \(q^*\mathcal G\) is flat over \(T\) and has point rank \(e\) at every
  point of \(T\).  Then \(q\) factors uniquely through the closed immersion
  of the rank-\(e\) stratum \(Z_e \hookrightarrow X\).
\end{theorem}

\begin{proof}
  \leanok
  \uses{thm:strata_ideal_le_ker, def:strata_data}
  The stratum \(Z_e\) is the closed subscheme cut out by the strata
  ideal-sheaf datum (\cref{def:strata_data}), and the kernel of its
  immersion is exactly that ideal sheaf.  A morphism into \(X\) factors
  through the closed subscheme cut out by a quasi-coherent ideal sheaf if
  and only if that ideal sheaf is contained in the kernel of the morphism,
  and then the factorization is unique because a closed immersion is a
  monomorphism.  By \cref{thm:strata_ideal_le_ker} the strata ideal sheaf is
  contained in the kernel of \(q\).
\end{proof}

\begin{theorem}[The rank-\(e\) stratum flattens \(\mathcal G\)]
  \label{thm:stratum_flattens}
  \lean{AlgebraicGeometry.Scheme.Modules.coherentSheafFlat_stratum}
  \leanok
  \uses{def:stratum, def:stratum_immersion, def:charts_cover,
    def:coherent_sheaf_flat, lem:pullback_isQuasicoherent_hom}
  \dcref{custom}
  Let \(\mathcal G\) be a quasi-coherent \(\OO_X\)-module admitting
  \(e\)-presentation charts, and \(\iota : Z_e \to X\) the closed immersion
  of the rank-\(e\) stratum.  Then \(\iota^*\mathcal G\) is flat over
  \(Z_e\).
\end{theorem}

\begin{proof}
  \leanok
  \uses{lem:presentation_pullback_sections, thm:strata_ideal_eq_entry_ideal,
    def:entry_ideal, lem:presentation_free_of_vanish,
    lem:sheaf_flat_of_flat_charts, def:strata_data}
  The preimages \(\iota^{-1}(V)\) of the \(e\)-presentation charts \(V\)
  form an affine open cover of \(Z_e\) (a closed immersion is an affine
  morphism, and the charts cover \(X\)).  Fix a chart \(V\) with
  presentation \(P = (\psi, \pi)\).  By
  \cref{lem:presentation_pullback_sections} the section module
  \(\Gamma(\iota^*\mathcal G, \iota^{-1}(V))\) has a matrix presentation
  with relation matrix \(\iota^\sharp(\psi)\).  Every entry of \(\psi\) lies
  in the entry ideal \(I(\psi)\) (\cref{def:entry_ideal}), which equals the
  strata ideal \(I_e(V)\) (\cref{thm:strata_ideal_eq_entry_ideal}); and the
  kernel of \(\iota^\sharp : \Gamma(X, V) \to \Gamma(Z_e, \iota^{-1}(V))\)
  is exactly \(I_e(V)\), the ideal cutting out the subscheme.  Hence
  \(\iota^\sharp(\psi) = 0\), so the section module is free
  (\cref{lem:presentation_free_of_vanish}), in particular flat, over
  \(\Gamma(Z_e, \iota^{-1}(V))\).  Conclude by
  \cref{lem:sheaf_flat_of_flat_charts}.
\end{proof}

\subsection{The chart locus}

\begin{theorem}[Chart supply on a locally noetherian scheme]
  \label{thm:presentation_chart_exists}
  \lean{AlgebraicGeometry.Scheme.Modules.exists_presentationChart_mem}
  \leanok
  \uses{def:presentation_chart, def:point_rank}
  \dcref{custom}
  Let \(S\) be a locally noetherian scheme and \(\F\) a finitely presented
  \(\OO_S\)-module.  Every point \(s \in S\) lies in an \(e\)-presentation
  chart for \(\F\) with \(e\) the point rank of \(\F\) at \(s\).
\end{theorem}

\begin{proof}
  \leanok
  \uses{lem:affine_finite_sections_nhds, thm:nakayama_prolongation,
    lem:point_rank_eq_chart_fiber_rank}
  Choose an affine open \(V \ni s\) on which \(M := \Gamma(\F, V)\) is a
  finite module over \(R := \Gamma(S, V)\)
  (\cref{lem:affine_finite_sections_nhds}).  Since \(S\) is locally
  noetherian, \(R\) is a noetherian ring, so the finite module \(M\) is
  finitely presented.  Let \(\mathfrak p\) be the prime of \(s\) in \(R\)
  and \(e\) the fiber rank of \(M\) at \(\mathfrak p\); by
  \cref{lem:point_rank_eq_chart_fiber_rank}, \(e\) is the point rank of
  \(\F\) at \(s\).  The Nakayama prolongation
  (\cref{thm:nakayama_prolongation}) supplies \(g \notin \mathfrak p\) such
  that every localization of \(M\) away from \(g\) admits a matrix
  presentation by \(e\) generators; by quasi-coherence
  \(\Gamma(\F, D(g))\) is such a localization, so the basic open
  \(D(g) \subseteq V\) is an \(e\)-presentation chart, and it contains
  \(s\) because \(g \notin \mathfrak p\).
\end{proof}

\begin{definition}[Chart locus]
  \label{def:chart_locus}
  \lean{AlgebraicGeometry.Scheme.Modules.chartLocus}
  \leanok
  \uses{def:presentation_chart}
  \dcref{custom}
  The \emph{\(e\)-chart locus} \(W_e \subseteq S\) of an \(\OO_S\)-module
  \(\F\) is the union of all \(e\)-presentation charts for \(\F\); it is an
  open subset of \(S\).
\end{definition}

\begin{lemma}[The chart locus contains the rank-\(e\) locus]
  \label{lem:chart_locus_contains_rank_locus}
  \dcref{custom}
  \lean{AlgebraicGeometry.Scheme.Modules.mem_chartLocus_of_pointRank_eq}
  \leanok
  \uses{def:chart_locus, def:point_rank}
  Let \(S\) be locally noetherian and \(\F\) finitely presented.  Every
  point of point rank \(e\) lies in \(W_e\).
\end{lemma}

\begin{proof}
  \leanok
  \uses{thm:presentation_chart_exists}
  \cref{thm:presentation_chart_exists} puts such a point inside an
  \(e\)-presentation chart.
\end{proof}

\begin{lemma}[Rank bound on the chart locus]
  \label{lem:chart_locus_rank_bound}
  \dcref{custom}
  \lean{AlgebraicGeometry.Scheme.Modules.pointRank_le_of_mem_chartLocus}
  \leanok
  \uses{def:chart_locus, def:point_rank}
  Let \(\F\) be a quasi-coherent \(\OO_S\)-module.  Every point of \(W_e\)
  has point rank at most \(e\).
\end{lemma}

\begin{proof}
  \leanok
  \uses{def:presentation_chart, def:fiber_rank,
    lem:point_rank_eq_chart_fiber_rank}
  A point \(s \in W_e\) lies in some \(e\)-presentation chart \(V\), so
  \(\Gamma(\F, V)\) is generated by \(e\) elements.  Its fiber
  \(\kappa(\mathfrak p_s) \otimes \Gamma(\F, V)\) at the prime of \(s\) is
  then spanned by \(e\) vectors, so the fiber rank at \(\mathfrak p_s\)
  (\cref{def:fiber_rank}) --- the point rank of \(\F\) at \(s\), by
  \cref{lem:point_rank_eq_chart_fiber_rank} --- is at most \(e\).
\end{proof}

\begin{lemma}[The restricted module is chart-covered]
  \label{lem:charts_cover_chart_locus}
  \dcref{custom}
  \lean{AlgebraicGeometry.Scheme.Modules.chartsCover_chartLocus}
  \leanok
  \uses{def:chart_locus, def:charts_cover, lem:pullback_isQuasicoherent_hom}
  Let \(\F\) be a quasi-coherent \(\OO_S\)-module.  The restriction of
  \(\F\) to the open subscheme \(W_e\) --- its pullback along the open
  immersion \(W_e \hookrightarrow S\) --- admits \(e\)-presentation charts.
\end{lemma}

\begin{proof}
  \leanok
  \uses{lem:presentation_pullback_sections, def:presentation_chart}
  Every point of \(W_e\) lies in some \(e\)-presentation chart
  \(V \subseteq W_e\).  The preimage of \(V\) in the open subscheme is an
  affine open mapped isomorphically onto \(V\), and by
  \cref{lem:presentation_pullback_sections} the section module of the
  restricted sheaf over it inherits a matrix presentation by \(e\)
  generators.
\end{proof}

\subsection{The rank stratum over the chart locus}

\begin{definition}[Rank-\(e\) stratum over the base]
  \label{def:rank_stratum}
  \lean{AlgebraicGeometry.Scheme.Modules.rankStratum}
  \leanok
  \uses{def:stratum, def:chart_locus, lem:charts_cover_chart_locus}
  \dcref{custom}
  The \emph{rank-\(e\) stratum} \(S_e\) of a quasi-coherent
  \(\OO_S\)-module \(\F\) is the rank-\(e\) stratum (\cref{def:stratum}) of
  the restriction of \(\F\) to the \(e\)-chart locus \(W_e\), which admits
  \(e\)-presentation charts by \cref{lem:charts_cover_chart_locus}: a
  closed subscheme of the open subscheme \(W_e\), hence a locally closed
  subscheme of \(S\).
\end{definition}

\begin{definition}[Immersion of the rank-\(e\) stratum]
  \label{def:rank_stratum_immersion}
  \lean{AlgebraicGeometry.Scheme.Modules.rankStratumι}
  \leanok
  \uses{def:rank_stratum, def:stratum_immersion, def:chart_locus}
  \dcref{custom}
  The canonical morphism \(\iota_e : S_e \to S\) is the composite of the
  closed immersion \(S_e \hookrightarrow W_e\)
  (\cref{def:stratum_immersion}) with the open immersion
  \(W_e \hookrightarrow S\); it is an immersion.
\end{definition}

\begin{theorem}[Support of the rank-\(e\) stratum over the base]
  \label{thm:mem_range_rank_stratum}
  \lean{AlgebraicGeometry.Scheme.Modules.mem_range_rankStratumι_iff}
  \leanok
  \uses{def:rank_stratum_immersion, def:point_rank}
  \dcref{custom}
  Let \(S\) be locally noetherian and \(\F\) finitely presented.  A point
  \(s \in S\) lies in the image of \(\iota_e : S_e \to S\) if and only if
  the point rank of \(\F\) at \(s\) equals \(e\).
\end{theorem}

\begin{proof}
  \leanok
  \uses{thm:mem_range_stratum, thm:point_rank_pullback,
    lem:chart_locus_contains_rank_locus}
  Write \(\F' := \F|_{W_e}\) for the restriction to the chart locus; by
  \cref{thm:point_rank_pullback} (applied to the open immersion) the point
  rank of \(\F'\) at a point of \(W_e\) equals the point rank of \(\F\) at
  its image in \(S\).

  If \(s = \iota_e(z)\), then \(s \in W_e\) and \(s\), viewed in \(W_e\),
  lies in the support of the stratum of \(\F'\), so the point rank of
  \(\F'\) there is \(e\) (\cref{thm:mem_range_stratum}); hence the point
  rank of \(\F\) at \(s\) is \(e\).  Conversely, if the point rank of
  \(\F\) at \(s\) is \(e\), then \(s \in W_e\)
  (\cref{lem:chart_locus_contains_rank_locus}); the point rank of \(\F'\)
  at \(s\) is \(e\), so \(s\) lies in the support of the stratum
  (\cref{thm:mem_range_stratum}), i.e.\ in the image of \(\iota_e\).
\end{proof}

\begin{theorem}[The rank-\(e\) stratum flattens \(\F\)]
  \label{thm:rank_stratum_flattens}
  \lean{AlgebraicGeometry.Scheme.Modules.coherentSheafFlat_rankStratum}
  \leanok
  \uses{def:rank_stratum, def:rank_stratum_immersion,
    def:coherent_sheaf_flat}
  \dcref{custom}
  Let \(\F\) be a quasi-coherent \(\OO_S\)-module.  The pullback
  \(\iota_e^*\F\) is flat over \(S_e\).
\end{theorem}

\begin{proof}
  \leanok
  \uses{thm:stratum_flattens, lem:charts_cover_chart_locus}
  Since \(\iota_e\) is the composite of the closed immersion
  \(\iota : S_e \to W_e\) with the open immersion \(W_e \to S\), pullback
  functoriality gives \(\iota_e^*\F \cong \iota^*(\F|_{W_e})\).  The
  restriction \(\F|_{W_e}\) admits \(e\)-presentation charts
  (\cref{lem:charts_cover_chart_locus}), so \(\iota^*(\F|_{W_e})\) is flat
  over \(S_e\) by \cref{thm:stratum_flattens}; flatness is invariant under
  isomorphism of modules.
\end{proof}

\subsection{Rank bound, open rank fibers, and the single-stratum
  factorization}

\begin{lemma}[The point rank is bounded on a noetherian scheme]
  \label{lem:point_rank_bounded}
  \lean{AlgebraicGeometry.Scheme.Modules.exists_pointRank_le}
  \leanok
  \uses{def:point_rank}
  \dcref{custom}
  Let \(S\) be a noetherian scheme and \(\F\) a finitely presented
  \(\OO_S\)-module.  There is an \(N\) such that the point rank of \(\F\)
  is at most \(N\) at every point of \(S\).
\end{lemma}

\begin{proof}
  \leanok
  \uses{thm:presentation_chart_exists, def:chart_locus,
    lem:chart_locus_rank_bound}
  For each \(s \in S\) choose an \(e(s)\)-presentation chart \(V_s \ni s\)
  with \(e(s)\) the point rank at \(s\)
  (\cref{thm:presentation_chart_exists}).  The noetherian space \(S\) is
  quasi-compact, so finitely many charts \(V_{s_1}, \dots, V_{s_k}\) cover
  \(S\); set \(N := \max_i e(s_i)\).  Any \(x \in S\) lies in some
  \(V_{s_i}\), hence in the \(e(s_i)\)-chart locus, so its point rank is at
  most \(e(s_i) \le N\) (\cref{lem:chart_locus_rank_bound}).
\end{proof}

\begin{theorem}[Openness of the rank fibers of a flat pullback]
  \label{thm:rank_fibers_open}
  \dcref{custom}
  \lean{AlgebraicGeometry.Scheme.Modules.isOpen_pointRank_pullback_eq}
  \leanok
  \uses{def:point_rank, def:coherent_sheaf_flat,
    lem:pullback_isQuasicoherent_hom}
  Let \(S\) be locally noetherian, \(\F\) a finitely presented
  \(\OO_S\)-module, and \(\varphi : T \to S\) a morphism with
  \(\varphi^*\F\) flat over \(T\).  For every \(e\), the locus of points of
  \(T\) at which the point rank of \(\varphi^*\F\) equals \(e\) is open in
  \(T\).
\end{theorem}

\begin{proof}
  \leanok
  \uses{lem:affine_finite_sections_nhds, lem:exists_matrix_presentation,
    lem:presentation_pullback_sections, lem:matrix_presentation_fp,
    lem:flat_sections_of_sheaf_flat, thm:stalk_rank_eq_point_rank}
  Let \(t_0\) be a point of the locus.  Choose an affine open
  \(V \ni \varphi(t_0)\) on which \(\Gamma(\F, V)\) is finite over the
  noetherian ring \(\Gamma(S, V)\)
  (\cref{lem:affine_finite_sections_nhds}), hence finitely presented, and
  an affine open \(W \ni t_0\) with \(W \subseteq \varphi^{-1}(V)\).  Set
  \(M := \Gamma(\varphi^*\F, W)\) and \(R := \Gamma(T, W)\).  Then \(M\) is
  flat over \(R\) (\cref{lem:flat_sections_of_sheaf_flat}) and finitely
  presented: a matrix presentation of \(\Gamma(\F, V)\)
  (\cref{lem:exists_matrix_presentation}) pulls back to a matrix
  presentation of \(M\) (\cref{lem:presentation_pullback_sections}), and a
  module with a matrix presentation is finitely presented
  (\cref{lem:matrix_presentation_fp}).

  A flat, finitely presented module is locally free, so the function
  \(\mathfrak p \mapsto \operatorname{rk} M_{\mathfrak p}\) is locally
  constant on the prime spectrum of \(R\).  Under the canonical
  identification of \(W\) with that spectrum, this function is the point
  rank of \(\varphi^*\F\) (\cref{thm:stalk_rank_eq_point_rank}).  Hence the
  set of points of \(W\) of point rank \(e\) is open in \(W\) and contains
  \(t_0\); as \(W\) is open in \(T\), the locus is a neighbourhood of each
  of its points.
\end{proof}

\begin{theorem}[Unique factorization through the rank-\(e\) stratum of the
  base]
  \label{thm:factor_rank_stratum_unique}
  \lean{AlgebraicGeometry.Scheme.Modules.existsUnique_factor_rankStratum}
  \leanok
  \uses{def:rank_stratum, def:rank_stratum_immersion,
    def:coherent_sheaf_flat, def:point_rank}
  \dcref{custom}
  Let \(S\) be locally noetherian, \(\F\) a finitely presented
  \(\OO_S\)-module, and \(\varphi : T \to S\) a morphism such that
  \(\varphi^*\F\) is flat over \(T\) and the point rank of \(\F\) at
  \(\varphi(t)\) equals \(e\) for every \(t \in T\).  Then \(\varphi\)
  factors uniquely through \(\iota_e : S_e \to S\).
\end{theorem}

\begin{proof}
  \leanok
  \uses{lem:chart_locus_contains_rank_locus, thm:point_rank_pullback,
    lem:charts_cover_chart_locus, thm:stratum_lift_unique}
  Every \(\varphi(t)\) has point rank \(e\), hence lies in the chart locus
  \(W_e\) (\cref{lem:chart_locus_contains_rank_locus}), so \(\varphi\)
  factors through the open immersion \(W_e \hookrightarrow S\) as
  \(\varphi = (W_e \hookrightarrow S) \circ q\) with \(q : T \to W_e\).
  Write \(\F' := \F|_{W_e}\); it admits \(e\)-presentation charts
  (\cref{lem:charts_cover_chart_locus}).  By functoriality of pullback,
  \(q^*\F' \cong \varphi^*\F\) is flat over \(T\), and by
  \cref{thm:point_rank_pullback} (applied twice, to \(q\) and to the open
  immersion) its point rank at \(t\) equals the point rank of \(\F\) at
  \(\varphi(t)\), namely \(e\).  Now \cref{thm:stratum_lift_unique} gives a
  unique \(l : T \to S_e\) with \(\iota \circ l = q\) for the closed
  immersion \(\iota : S_e \to W_e\); composing,
  \(\iota_e \circ l = \varphi\).

  For uniqueness, let \(l' : T \to S_e\) satisfy
  \(\iota_e \circ l' = \varphi\).  Then \(\iota \circ l'\) and \(q\) are two
  factorizations of \(\varphi\) through the open immersion
  \(W_e \hookrightarrow S\), which is a monomorphism, so
  \(\iota \circ l' = q\); by the uniqueness in
  \cref{thm:stratum_lift_unique}, \(l' = l\).
\end{proof}

\section{Universal property of the strata}
\label{sec:flatstrat_universal}

\begin{theorem}[Universal property of the flat-locus stratification (case
  \(n = 0\))]
  \label{thm:flattening_stratification_universal}
  \lean{AlgebraicGeometry.flatLocusStratification_universal}
  \leanok
  \uses{def:coherent_sheaf_flat}
  \dcref{custom}
  Let \(S\) be a noetherian scheme and \(\F\) a coherent \(\OO_S\)-module.
  Then there exist a finite index set \(I\) and locally-closed immersions
  \(\iota_f : S_f \to S\) (\(f \in I\)) with pairwise disjoint images
  set-theoretically covering \(|S|\), such that the pullback of \(\F\)
  along \(\iota_f\) is flat over \(S_f\) for every \(f\) (in the sense of
  \cref{def:coherent_sheaf_flat}), and such that, writing
  \(S' := \coprod_{f \in I} S_f\) and \(i : S' \to S\) for the morphism
  assembled from the \(\iota_f\), the pair \((S', i)\) has the universal
  property: every morphism \(\varphi : T \to S\) for which
  \(\varphi^* \F\) is flat over \(T\) factors \emph{uniquely} through
  \(i\).

  For a finitely presented module, flatness is locally equivalent to local
  freeness, and the rank decomposes \(T\) into clopen pieces.  Conversely,
  every morphism factoring through \(i\) has flat pullback by stability of
  flatness under base change (\cref{lem:sheaf_flat_pullback}).  For a relative
  projective family, the analogous universal property additionally uses the
  cohomology-and-base-change layer of
  \cref{lem:noetherian_induction_strata}.  Reduced strata alone do not give
  such a universal property for an arbitrary proper family: over
  \(S = \Spec k[\varepsilon]\) with the flat module \(\F = \OO_S\), the
  identity must factor through the flattening stratum, which forces its
  canonical non-reduced scheme structure.
\end{theorem}

\begin{proof}
  \leanok
  \uses{lem:point_rank_bounded, def:rank_stratum, def:rank_stratum_immersion,
    thm:mem_range_rank_stratum, thm:rank_stratum_flattens,
    lem:sheaf_flat_pullback, thm:rank_fibers_open, thm:point_rank_pullback,
    thm:factor_rank_stratum_unique, def:point_rank}
  \textbf{The strata.}  By \cref{lem:point_rank_bounded} there is an \(N\)
  bounding the point rank of \(\F\) on \(S\).  Take
  \(I := \{0, 1, \dots, N\}\) and, for \(e \in I\), the rank-\(e\) stratum
  \(\iota_e : S_e \to S\)
  (\cref{def:rank_stratum} and \cref{def:rank_stratum_immersion}); each \(\iota_e\) is
  an immersion.  By \cref{thm:mem_range_rank_stratum} the image of
  \(\iota_e\) is exactly the locus of point rank \(e\); every point of
  \(S\) has a point rank, which is at most \(N\), so the images cover
  \(|S|\), and they are pairwise disjoint because the point rank is a
  single natural number.  Flatness of \(\iota_e^*\F\) over \(S_e\) is
  \cref{thm:rank_stratum_flattens}.

  \textbf{Factorization.}  Let \(\varphi : T \to S\) have \(\varphi^*\F\)
  flat over \(T\).  For \(e \in I\) let \(T_e \subseteq T\) be the locus
  where the point rank of \(\varphi^*\F\) equals \(e\).  By
  \cref{thm:point_rank_pullback} the point rank of \(\varphi^*\F\) at \(t\)
  is the point rank of \(\F\) at \(\varphi(t)\), hence at most \(N\); so
  the \(T_e\) cover \(T\), they are pairwise disjoint, and each is open by
  \cref{thm:rank_fibers_open}.  Thus \(T\) is the disjoint union of the
  open subschemes \(T_e\), and a morphism out of \(T\) is the same as a
  family of morphisms out of the \(T_e\).

  On a piece, the composite
  \(T_e \hookrightarrow T \xrightarrow{\varphi} S\) has flat pullback of
  \(\F\) --- it is the restriction of \(\varphi^*\F\) to the open subscheme
  \(T_e\), flat by \cref{lem:sheaf_flat_pullback} --- and, by
  \cref{thm:point_rank_pullback} again, the point rank of \(\F\) at the
  image of any \(t \in T_e\) is \(e\).  By
  \cref{thm:factor_rank_stratum_unique} there is a unique
  \(l_e : T_e \to S_e\) with \(\iota_e \circ l_e = \varphi|_{T_e}\).
  Define \(\psi : T \to \coprod_{f \in I} S_f\) on \(T_e\) as \(l_e\)
  followed by the coprojection of \(S_e\); then \(i \circ \psi = \varphi\),
  because this holds on each piece of the cover.

  \textbf{Uniqueness.}  Let \(\psi' : T \to \coprod_f S_f\) also satisfy
  \(i \circ \psi' = \varphi\).  Fix \(e\) and \(t \in T_e\), and let
  \(e'\) be the index of the summand containing \(\psi'(t)\), say
  \(\psi'(t) = z \in S_{e'}\).  Then
  \(\varphi(t) = i(\psi'(t)) = \iota_{e'}(z)\) lies in the image of
  \(\iota_{e'}\), so the point rank of \(\F\) at \(\varphi(t)\) is \(e'\)
  (\cref{thm:mem_range_rank_stratum}); it is also \(e\), since
  \(t \in T_e\) (\cref{thm:point_rank_pullback}).  Hence \(e' = e\): the
  morphism \(\psi'\) maps \(T_e\) into the open summand \(S_e\) of the
  coproduct, so \(\psi'|_{T_e}\) factors through the coprojection of
  \(S_e\) as some \(l'_e : T_e \to S_e\), and
  \(\iota_e \circ l'_e = i \circ \psi'|_{T_e} = \varphi|_{T_e}\).  The
  uniqueness in \cref{thm:factor_rank_stratum_unique} forces
  \(l'_e = l_e\); therefore \(\psi'\) and \(\psi\) agree on every piece
  \(T_e\), and \(\psi' = \psi\).
\end{proof}

\section{Application: relative curves}
\label{sec:flatstrat_curve_application}

Let \(C\) be a smooth proper curve over a field \(k\) and \(T\) a noetherian
\(k\)-scheme.  To apply \cref{thm:flattening_stratification_exists} to
\(C \times_k T \to T\), it remains to pass from properness to projectivity.

\begin{lemma}
[Smooth proper curve over \(k\) is projective]
  \label{lem:smooth_proper_curve_projective}
  \dcref{custom}
  \uses{def:projective_morphism}
  \textit{Source: classical; cf.\ \cite{hartshorne-ag}, II.6.7 \& II.7.6 (curves).}
  Let \(C\) be a smooth proper curve over a field \(k\).  Then \(C \to \Spec k\) is
  projective.  Consequently, for any noetherian \(k\)-scheme \(T\), the relative
  curve \(C \times_k T \to T\) is projective in the sense of Grothendieck (there
  exists a relatively very ample line bundle on \(C \times_k T\) over \(T\)).
\end{lemma}

\begin{proof}
  Projectivity of a smooth proper curve over a field is classical
  (\cite{hartshorne-ag}, II.6.7 / II.7.6; alternatively: every divisor of large
  positive degree on a smooth proper curve \(C\) is very ample, providing the
  embedding \(C \hookrightarrow \PP^N_k\) for some \(N\)). Pulling the ample line
  bundle back along \(C \times_k T \to C\) produces a relatively very ample line
  bundle on \(C \times_k T\) over \(T\).
\end{proof}

\begin{corollary}

[Flattening stratification for a coherent sheaf on a relative curve]
  \label{cor:flattening_stratification_curve}
  \dcref{custom}
  \uses{lem:smooth_proper_curve_projective, thm:flattening_stratification_exists}
  Let \(k\) be a field, \(C\) a smooth proper curve over \(k\), and \(T\) a noetherian
  \(k\)-scheme.  Let \(\F\) be a coherent sheaf on \(C \times_k T\). Then the
  flattening-stratification conclusion of
  \cref{thm:flattening_stratification_exists} applies: there is a finite
  set \(I\) of Hilbert polynomials and locally-closed subschemes
  \(T_f \hookrightarrow T\) for \(f \in I\), set-theoretically covering \(T\)
  disjointly, such that (a) \(\F|_{C \times_k T_f}\) is flat over \(T_f\) for
  each \(f\), and (b), after setting \(T'':=\coprod_{f\in I}T_f\), any
  morphism \(\varphi : T' \to T\) such that \(\varphi^* \F\) on
  \(C \times_k T'\) is \(T'\)-flat factors uniquely through the canonical
  morphism \(T'' \to T\).
\end{corollary}

\begin{proof}
  By \cref{lem:smooth_proper_curve_projective}, \(C \times_k T \to T\) is
  projective, hence factors through a closed immersion into \(\PP^n_T\) for some
  \(n \ge 0\).  Push \(\F\) forward along the closed immersion (extending by zero
  outside the image) to a coherent sheaf \(\widetilde{\F}\) on \(\PP^n_T\); this
  operation preserves \(T\)-flatness in both directions because the closed
  immersion is a finite morphism.  Apply
  \cref{thm:flattening_stratification_exists} to \(\widetilde{\F}\) on
  \(\PP^n_T\) to obtain \(\{T_f\}\) with the stated universal property.
\end{proof}

\section{Mathematical ingredients}
\label{sec:flatstrat_mathlib_status}

The construction combines four kinds of input.  First, module flatness is
stable under localization and arbitrary base change, and flatness of a
quasi-coherent module over a base may be checked on affine charts.  Second,
generic freeness over a noetherian domain supplies the dense flat open used
in the Noetherian induction of \cref{lem:flat_locus_reduction_lean}.  Third,
entry ideals of finite presentations define the rank strata and their
universal property, as developed in \cref{sec:entry_ideal}--
\cref{sec:flatstrat_representability}.  Finally, for a projective family,
Serre vanishing and cohomology and base change compare the fibre Hilbert
polynomial with the ranks of the finitely many direct images
\(\pi_*\F(N),\ldots,\pi_*\F(N+n)\).  These inputs produce the canonical
Hilbert-polynomial indexing in
\cref{thm:flattening_stratification_exists}.

\section{Relation to adjacent constructions}
\label{sec:flatstrat_out_of_scope}

Castelnuovo--Mumford regularity supplies the uniform twists used in the
projective cohomology argument, while flattening stratification supplies the
locally closed condition in the construction of Quot schemes: a family of
quotients has prescribed Hilbert polynomial precisely after its classifying
morphism factors through the corresponding stratum.  For a smooth proper
curve, \cref{cor:flattening_stratification_curve} gives this input after base
change to an arbitrary noetherian scheme.  Passing from Quot schemes to the
representability of the relative Picard functor is the quotient route of
\cref{sec:fga_pic_setup}, a separate argument resting on a descent step the
formalisation instead avoids by the Milne--Koll\'ar route; either way it does
not alter the flattening construction proved here.
